%% Compile: lualatex x2 (canonical engine since the 2026-06-28 R-TikZ figure
%% migration). The PeerJ upload workspace uses pdflatex instead -- see
%% papers/VERSIONS.md.
\documentclass[11pt]{article}

\usepackage[margin=1in]{geometry}
\usepackage{graphicx}
\usepackage{amsmath,amssymb}
\usepackage{amsthm}
% defaultfeatures={Renderer=Node}: newtxtext forces Renderer=Basic under LuaTeX,
% which drops the TeX ``/'' ligatures (quotes rendered as two single-quote glyphs);
% the Node renderer restores them.
\usepackage[defaultfeatures={Renderer=Node}]{newtxtext}
\usepackage{newtxmath}
\usepackage{booktabs}
\usepackage{array}
\usepackage{float}
\usepackage{microtype}
\usepackage[round]{natbib}
\usepackage{hyperref}
\usepackage{threeparttable}
\usepackage{orcidlink}
\usepackage{tcolorbox}
\tcbuselibrary{skins,breakable}
% Muted dark link colors instead of boxed links (de-noises Table 1 Study cells)
\definecolor{linkblue}{RGB}{13,51,115}
\definecolor{citegreen}{RGB}{10,77,56}
\definecolor{urlplum}{RGB}{77,26,90}
\hypersetup{
colorlinks=true, linkcolor=linkblue, citecolor=citegreen, urlcolor=urlplum,
  pdftitle={Apparent emergence timing in small-scale grokking testbeds can be scan-statistic false positives: exact rates and a sustained-crossing correction},
  pdfauthor={Zeyu Fu},
  pdfsubject={Scan-statistic calibration of apparent emergence timing in small-scale grokking testbeds},
  pdfkeywords={emergence, grokking, scan statistics, false-positive rate, sustained crossing, calibration}
}

\theoremstyle{plain}
\newtheorem{theorem}{Theorem}
\newtheorem{lemma}{Lemma}
\newtheorem{proposition}{Proposition}
\newtheorem{corollary}{Corollary}
\theoremstyle{definition}
\newtheorem{remark}{Remark}

\newenvironment{mathderivation}[1]{%
  \begin{tcolorbox}[
    colback=gray!6, colframe=gray!40, arc=2pt,
    left=6pt, right=6pt, top=4pt, bottom=4pt,
    title={\small\bfseries #1},
    fonttitle=\small\bfseries,
    breakable, enhanced
  ]%
}{%
  \end{tcolorbox}%
}
\input{../figs/figpreamble.tex}

% Tighter float placement to reduce whitespace
\renewcommand{\topfraction}{0.9}
\renewcommand{\bottomfraction}{0.9}
\renewcommand{\textfraction}{0.07}
\renewcommand{\floatpagefraction}{0.75}
\setcounter{topnumber}{3}
\setcounter{bottomnumber}{2}
\setcounter{totalnumber}{5}

\begin{document}

\title{Apparent emergence timing in small-scale grokking testbeds can be
scan-statistic false positives: exact rates and a sustained-crossing correction}
\author{Zeyu Fu\,\orcidlink{0009-0001-8329-0108}\\
State Key Laboratory of Trauma and Chemical Poisoning, Institute of Combined Injury,\\
Chongqing Engineering Research Center for Nanomedicine, College of Preventive Medicine,\\
Army Medical University, Chongqing, China\\
\href{mailto:fuzeyu09@gmail.com}{fuzeyu09@gmail.com}}
\date{}

\maketitle

\begin{abstract}
Threshold-crossing ``emergence detectors'' can manufacture apparent phase transitions out
of evaluation noise. Re-analysing the public BIG-Bench data behind the original
emergent-abilities claims, an exact scan-statistics computation shows that a
chance$+0.10$ detector is expected to manufacture $17.7$ false ``emergent tasks'' out of
$89$ scanned; a look-elsewhere correction leaves most observed detections standing while
flagging several as unpowered. Crucially, the flagged firings are detector artifacts on
the public score release, not claims made in the source paper, so the audit corroborates
the published claims rather than debunking them. This failure mode reduces to classical
run/scan statistics. Under a no-transition null with per-checkpoint exceedance
probability $q$, the single-crossing false-positive rate $1-(1-q)^{M}$ rises with the
number of checkpoints $M$, so refining the evaluation grid strictly worsens it.
Requiring $K$ consecutive exceedances suppresses false positives; the exact Markov-chain
suppression factor at the evaluated operating point ($q=0.00272$, $K=2$) is $286\times$.
We give the exact finite-Markov-chain run-length probability, a
power/false-negative/latency tradeoff, and a closed-form nomogram that picks $K$ from
$M$, $q$, and a false-positive budget $\alpha$. On an in-context temporal-difference
testbed with pre-registered threshold and criteria, the single-crossing detector fires
in $42/45$ non-learning runs at $\tau=0.7$; a per-run mixture model reproduces the entire
sustained-$K$ ladder to within $1.5$ runs, and its prescribed sustained-$5$ criterion
reduces the count to $1/45$---an observed ${\sim}42\times$ reduction. The account extends
to drifted baselines, threshold and evaluation-set-size ladders, public Pythia
checkpoints, and three apparent effects. Apparent small-scale emergence timing therefore
needs no new phenomenon: it can often be scan-statistic false positives,
correctable by a stated calibration checklist.
\end{abstract}

\vspace{4pt}\noindent\textbf{Keywords:} emergence; grokking; false-positive rate; scan
statistics; run statistics; sustained-crossing criterion; calibration; in-context learning;
temporal-difference learning; small-scale testbeds

% Highlights (Elsevier highlights go in a separate file at submission):
% - At the pre-registered tau=0.7 bar, a single-crossing detector fires in 42/45 non-learning runs.
% - The mixture-prescribed sustained-5 criterion reduces this to 1/45 (~42x observed reduction).
% - "Emergence rate rises with horizon" is measured baseline drift, not learning.
% - Exact run statistics yield a preflight nomogram for choosing K from (N, p, tau, M).
% - A solved-task positive control is a necessary preflight for emergence claims.

\section{Introduction}\label{sec:intro}

Phase-transition-like learning---where a model's capability appears to jump
discontinuously as training or scale proceeds---has attracted substantial empirical
and theoretical attention since the characterisation of ``emergent abilities'' by
\citet{wei2022emergent} and the grokking phenomenon by \citet{power2022grokking}.
Tiny-transformer testbeds make such transitions reproducible
at low cost, enabling systematic ablations over optimizers, data curriculum, and
architectural choices that would be prohibitive at scale.
Yet cheapness creates its own risk: the same experimenter who can run thousands of
cells can just as easily over-read them.
\citet{schaeffer2023mirage} demonstrated that apparent emergent abilities in large
language models can be artefacts of metric choice rather than genuine phase transitions;
related work on inverse scaling \citep{mckenzie2023inversescaling} has shown that
apparently robust scaling trends can reverse or disappear under closer scrutiny.
These findings motivate an analogous discipline for small-scale emergence work,
where the confounds are different (noise detectors, unlearnable targets, nuisance axes)
but the corrective logic is the same (Figure~\ref{fig:scheme}): verify that the measurement protocol can
distinguish a real signal from a plausible null before interpreting the result.

We make this discipline precise for the small-scale regime by reducing the
emergence-timing readout itself to a classical statistical problem. A threshold-crossing
``emergence detector'' that scans a scalar score across $M$ evaluation checkpoints and
fires on the first crossing is a scan statistic; under a chance-accuracy null its
false-positive behaviour is governed by run and scan statistics
\citep{feller1968,glaz2001scan,fulou2003runs}. The detector is not a straw man: it is how
emergence is operationalized across the literature---performance clearing an above-random
bar at some scale point \citep{wei2022emergent}, the BIG-Bench ``breakthroughness''
metric \citep{srivastava2022bigbench}, an ability manifesting once the pretraining loss
crosses a threshold \citep{du2024loss}, the emergence point predicted from pre-emergence
checkpoints \citep{snell2024predicting}, and, in the grokking and induction-head
literatures, onset timing read as the first checkpoint at which a score clears a fixed
threshold \citep{power2022grokking,olsson2022induction}. Section~\ref{sec:theorem} shows that the
noise-only single-crossing rate is $1-(1-q)^{M}$ and rises toward one as the evaluation
grid is refined, that a sustained-$K$ criterion suppresses it by $q^{-(K-1)}$, and that a
simple nomogram (Section~\ref{sec:nomogram}) calibrates $K$ to a target false-positive
budget. We then validate this account to within $2\%$ on a real in-context
temporal-difference (TD) testbed (Section~\ref{sec:r2}) and show that three apparent
effects---an ordering timing effect, an ``emergence rises with horizon'' scaling, and
early-trajectory predictiveness---reduce to this false-positive mechanism or to a companion
nuisance-axis confound (Sections~\ref{sec:r1},~\ref{sec:r3}). The framing is a
reconciliation rather than a new phenomenon, in the genre of \citet{oneto2026}, who
reconcile grokking with existing statistical learning theory without positing new
mechanisms; here the known mechanism is scan-statistic false positives. The mathematics we
use is classical; the contribution is the reduction and a calibrated instrument for
small-scale emergence science.

The in-context-emergence literature is almost entirely composed of positive claims
(Figure~\ref{fig:landscape}, Table~\ref{tab:emergence-landscape}); what it lacks is a calibrated negative that distinguishes
genuine emergence from a detector firing on noise---precisely what the reduction supplies.

\begin{figure}[t]
\centering
\figtikz[\linewidth]{E2_scheme}
\caption{Study design and headline finding, drawn at the pre-registered $\tau=0.7$ bar. On a stationary evaluation trace (an in-context TD testbed that does not acquire the task within the studied budget), a single-threshold-crossing emergence detector fires in $42/45$ runs, exactly as classical run statistics predict; requiring $K$ consecutive exceedances, with $K=5$ read from the per-run mixture of Section~\ref{sec:theorem}, reduces this to $1/45$ against $1.5$ predicted. At the illustrative $\tau=0.8$ bar used for the case-study figure the same chain runs $26/45\to0/45$ at $K=2$. The trace sketch is schematic: at $\tau=0.7$ a chance-level run crosses the bar repeatedly in isolation, and it is the absence of a long run, not the absence of crossings, that the criterion tests. The criterion is validated in three settings: a mildly drifted baseline ($0/45$ at $\tau=0.8$; the grid plateaus $0.17$ above its constant-predictor floor), a strongly drifted one ($0.35$ above floor, where the nomogram-corrected $K$ restores calibration), and genuine transitions ($90/90$ induction emergences detected at zero added latency).}
\label{fig:scheme}
\end{figure}

\begin{table}[htbp]
\centering
\caption{Published in-context-emergence claims vs this work's calibrated outcomes---an
expanded version of the landscape summarised in Figure~\ref{fig:landscape}.}
\label{tab:emergence-landscape}
\small
\begin{tabular}{@{}>{\raggedright\arraybackslash}p{0.18\linewidth}>{\raggedright\arraybackslash}p{0.17\linewidth}>{\raggedright\arraybackslash}p{0.24\linewidth}>{\raggedright\arraybackslash}p{0.26\linewidth}@{}}
\toprule
Study & Phenomenon & Claim & Key figure \\
\midrule
\citet{olsson2022induction} & induction / copy & emerges (phase change) & $2.5$--$5{\times}10^{9}$ tokens \\
\citet{lee2023icrl} & in-context reinforcement learning (RL) & emerges (pretraining) & near-optimal \\
\citet{laskin2022ad} & in-context RL & emerges (distillation) & improves policy in-context \\
\citet{zisman2024noise} & in-context RL & emerges (noise distill.) & Dark Room $6.36\!\to\!14.08$ \\
\citet{wang2024tdtransformers} & in-context TD & emerges (forward pass) & value error $\downarrow$ with context \\
\citet{schaeffer2023mirage} & (critique) & emergence may be a metric artifact & --- \\
\citet{sabry2025induction} & (critique) & induction signature $\neq$ load-bearing & anti-induction ${\approx}0$ \\
\textbf{This work} & in-context TD & calibration-negative & $42/45\!\to\!1/45$ ($\tau{=}0.7$, $K{=}5$); $26/45\!\to\!0/45$ ($\tau{=}0.8$, $K{=}2$); plateau $0.562$, floor $0.390$ \\
\bottomrule
\end{tabular}
\end{table}

\subsection{Related work}\label{sec:related}
\begin{figure}[t]
\centering
\figtikz[0.88\linewidth]{E2_landscape}
\caption{Published in-context-emergence claims versus this
work's calibrated outcomes. Induction/copy \citep{olsson2022induction}, in-context RL
\citep{lee2023icrl}, and in-context TD (Wang 2024; \citealp{wang2024tdtransformers}) are each reported
to emerge, while \citet{schaeffer2023mirage} caution that emergence can be a
metric/measurement artifact. We confirm a clean positive for induction
(${\approx}1850$ steps) but show the in-context-TD case is a calibration failure---the
chain annotated inside the panel is $42/45\to1/45$ ($\tau=0.7$): a single-crossing
detector manufactures emergence in $42/45$ runs that the mixture-prescribed sustained-$5$
criterion reduces to $1/45$---supplying the calibrated negative the positive-claim
literature lacked. The $\tau=0.8$ single-crossing chain $26/45\to0/45$ (where $K=2$
already suffices) lives in Table~\ref{tab:tdgrid} and Figure~\ref{fig:td-grid}, not in
this panel.}
\label{fig:landscape}
\end{figure}


\paragraph{Emergence in large language models and whether it is a mirage.}
\citet{wei2022emergent} introduced the term ``emergent abilities'' for capabilities
that appear sharp and unpredicted by extrapolating smaller-model performance.
\citet{schaeffer2023mirage} challenged this framing: by switching from discontinuous
to continuous metrics, the apparent phase transitions smooth out, suggesting that
emergence is partially an artefact of metric nonlinearity rather than a genuine
discontinuity in the underlying computation.
\citet{mckenzie2023inversescaling} showed that some capabilities actually
decrease with scale, further complicating the picture of what ``emergence''
means empirically.
Our work brings the same measurement discipline to the small-scale
regime: just as metric choice determines whether large-model abilities look emergent,
detector choice determines whether small-model learning looks like a phase transition.
\citet{oneto2026} is a recent instance of the reconciliation genre we adopt---explaining a
grokking-type phenomenon through existing theory rather than a new mechanism.

\paragraph{Grokking as a testbed for emergence.}
\citet{power2022grokking} introduced the grokking phenomenon---delayed generalisation
long after overfitting---on small algorithmic datasets, making it a natural controlled
testbed for studying phase-transition-like learning.
Subsequent work has characterised the geometry \citep{liu2023omnigrok}, lazy-to-rich
transition \citep{kumar2024lazytorich}, norm
dynamics \citep{prieto2025stablemax}, and hidden-progress structure
\citep{nanda2023progress,barak2022hidden} of the grokking transition.
Early-trajectory prediction of grokking outcomes has been attempted by
\citet{notsawo2023predicting} and \citet{xu2026earlywarning}; our third case study
(Section~\ref{sec:r3}) shows that such predictions must freeze the optimizer axis before
the predictiveness can be attributed to run-level dynamics rather than configuration
identity.

\paragraph{In-context learning: circuits and emergence.}
The mechanistic study of in-context learning has centred on induction heads
\citep{olsson2022induction} as the circuit-level substrate for in-context
generalisation, including their statistical form on Markov data
\citep{edelman2024statinduction}.
Our ordering case study (Section~\ref{sec:r1}) exploits the sharp emergence of the
induction circuit as a measurable delay, testing whether data-presentation order can
compress it.
The broader in-context learning literature has examined how transformers implement
algorithm-like computations in-context \citep{akyurek2024iclalgorithms,garg2022icl,vonoswald2023gd}, but whether emergence timing is modifiable by curriculum has had little
direct experimental study at small scale.

\paragraph{In-context reinforcement learning (RL).}
\citet{lee2023icrl} showed that transformers pretrained with supervised imitation
on diverse RL trajectories can perform near-optimal in-context reinforcement learning
at test time, demonstrating that in-context RL emergence is achievable
in principle.
More directly, \citet{wang2024tdtransformers} report that transformers trained for
policy-evaluation tasks can discover and implement temporal-difference learning in their
forward pass, with value error decreasing as the in-context horizon grows; our TD case
study (Section~\ref{sec:r2}) is a calibration boundary for online single-task
in-context-TD---a harder regime than the pretrained multi-task setting of
\citet{wang2024tdtransformers} and \citet{lee2023icrl}---establishing the solved-reference
bar such a claim must clear before it is adjudicable on a single-task online testbed.
In our setup the same transformer family cannot be straightforwardly
trained to perform in-context temporal-difference learning above chance, and the naive
threshold detector fires on noise rather than on genuine learning.
This calibration failure is not a claim that in-context TD is impossible, but that
its testbed must provide a solved-reference cell before the detector can distinguish
emergence from noise.

\paragraph{Curriculum learning and data ordering.}
The idea that training on examples in order of increasing difficulty can improve
sample efficiency dates at least to \citet{bengio2009curriculum}, and is surveyed by \citet{soviany2022curriculum}.
Whether curriculum ordering helps with emergence specifically---compressing the
delay of a phase transition rather than improving asymptotic performance---is less
well studied.
Our ordering case study (Section~\ref{sec:r1}) provides a direct negative: on the
in-context induction task the ordering of training examples does not shorten the emergence
delay, and easy$\to$hard ordering mildly lengthens it.

\paragraph{Run and scan statistics.}
The false-positive analysis of Section~\ref{sec:theorem} is an application of classical
run and scan statistics. The distribution of success runs in Bernoulli sequences is
treated in \citet{feller1968}; the finite-Markov-chain-imbedding technique for exact
run-length distributions is developed by \citet{fulou2003runs}; and the scan-statistic
viewpoint on threshold exceedances over an ordered index is surveyed by
\mbox{\citet{glaz2001scan}}. Our use of these tools is standard; the novelty is their application
to emergence-timing detectors and the resulting calibrated instrument.

\paragraph{Early-training prediction and reproducibility.}
Predicting final performance from early training trajectories has attracted attention
as a cheap proxy for hyperparameter search \citep{notsawo2023predicting,xu2026earlywarning}.
Reproducibility and rigorous null-hypothesis discipline in machine learning more broadly
have been emphasised by \citet{pineau2021reproducibility}, who argue that negative
results and pre-registration are essential to combating the publication bias toward
positive findings---a discipline paralleled by the statistical-precipice critique of
point-estimate reporting in deep RL \citep{agarwal2021precipice}; we treat pre-registration as a methods virtue below rather than as the
identity of the paper.

\section{The false-positive theorem}\label{sec:theorem}

We formalise why a threshold-crossing detector manufactures apparent emergence, and why a
sustained criterion suppresses it. The underlying mathematics is classical run/scan
statistics \citep{feller1968,fulou2003runs,glaz2001scan}; what is new here is the reduction
of small-scale emergence-timing readouts to this setting, the power/latency analysis of the
sustained criterion (Lemma~\ref{lem:sustain}(c)), and the calibrated instrument of
Section~\ref{sec:nomogram}.

\paragraph{Setup.}
A detector observes a scalar signal $\hat p_1,\dots,\hat p_M$ over $M$ evaluation
checkpoints and declares ``emergence'' when the signal crosses a fixed threshold $\tau$.
Under the null---no real transition, the model at a stationary baseline accuracy $p$ over
$N$ evaluation
samples---each checkpoint reports $\hat p_t = \tfrac1N\sum_{i=1}^N\mathrm{Bernoulli}(p)$,
so $N\hat p_t\sim\mathrm{Binomial}(N,p)$ with marginal per-checkpoint exceedance
probability
\[
  q(\tau) = P\!\bigl(\mathrm{Binom}(N,p) \geq \lceil \tau N\rceil\bigr).
\]
What the null asserts is stationarity, not chance. The parameter $p$ is the level the
trace sits at while nothing is happening, which may be the task's chance level, a
partially learned plateau, or a baseline that drifts as some easier component of the task
is acquired. This distinction matters in practice, because every testbed analysed below
plateaus above its own trivial floor, and it is the plateau rather than the floor that
sets $q$.

\begin{theorem}[Single-crossing false-positive rate and its $M$-monotonicity]\label{thm:single}
Let the detector fire if $\hat p_t \geq \tau$ for at least one $t\in\{1,\dots,M\}$.
\begin{enumerate}
\item[\textnormal{(a)}] If the $\hat p_t$ are independent, the single-crossing
  false-positive rate is exactly
  \[
    P_{\mathrm{single}}(M) = 1-(1-q)^{M},
  \]
  which is strictly increasing in $M$, satisfies $P_{\mathrm{single}}(M)\leq Mq$, and
  tends to $1$ as $M\to\infty$.
\item[\textnormal{(b)}] If the exceedance events $\{\hat p_t\geq\tau\}$ are positively
  associated in the FKG/Harris sense, then $P_{\mathrm{single}}(M)\leq 1-(1-q)^{M}$; the
  independence formula is an upper bound.
\end{enumerate}
\end{theorem}

\begin{proof}
(a)~Under independence the non-exceedance events each have probability $1-q$, so
$P(\text{no crossing})=(1-q)^{M}$ and $P_{\mathrm{single}}=1-(1-q)^{M}$. Strict
monotonicity in $M$ follows from $0<q<1$; the union bound gives $P_{\mathrm{single}}\leq
Mq$; and $(1-q)^{M}\to0$ gives the limit.
(b)~Positive association (the FKG/Harris inequality) implies that the decreasing events
$\{\hat p_t<\tau\}$ satisfy $P\!\bigl(\bigcap_{t}\{\hat p_t<\tau\}\bigr)\geq\prod_{t}
P(\hat p_t<\tau)=(1-q)^{M}$, whence $P_{\mathrm{single}}=1-P(\text{no crossing})\leq
1-(1-q)^{M}$.
\end{proof}

\begin{remark}[Association is a hypothesis, and it flips the bound direction]\label{rem:assoc}
Part~(b) is stated under an explicit positive-association hypothesis, not proved from
the model. Positive dependence reduces the number of effectively independent trials, so for
the single-crossing detector it makes the i.i.d.\ formula an upper bound. The
practical reading of part~(a) is that with $M=O(1/q)$ checkpoints the noise-only crossing
probability is already of order one, so finer evaluation grids strictly worsen the
false-positive rate. Empirically, the hypothesis is at most weakly active in our testbed:
the lag-1 autocorrelation of the mean-centered (linearly detrended) validation signal has
median $+0.04$ ($+0.03$) across the 45 in-context-TD logs, and $41/45$ detrended values lie
inside the $\pm1.96/\sqrt{M}=\pm0.138$ white-noise band. Because the run-length formulas
at $K\geq3$ are sensitive to dependence beyond one step, the audit extends to lags
$2$--$5$: the detrended medians are $+0.01$, $+0.01$, $-0.02$, and $-0.02$, with
$40$--$44$ of $45$ runs inside the same band at every lag. The marginal-variance match
reported below is therefore accompanied by measured near-independence out to lag 5 rather
than assumed independence.
\end{remark}

\begin{lemma}[Sustained-$K$ suppression and the power/latency tradeoff]\label{lem:sustain}
Require the threshold to be held for $K$ consecutive checkpoints.
\begin{enumerate}
\item[\textnormal{(a)}] Under independence, the sustained-$K$ false-positive rate
  $P_{\mathrm{sustain}\text{-}K}$ equals the probability of a success run of length $\geq K$
  in $M$ Bernoulli$(q)$ trials, given exactly by the finite-Markov-chain-imbedding
  recurrence on states $\{0,1,\dots,K\}$ (current consecutive-success run length, with $K$
  absorbing), with small-$q$ asymptotic
  \[
    P_{\mathrm{sustain}\text{-}K} = (M-K+1)\,q^{K}\,(1+o(1)).
  \]
\item[\textnormal{(b)}] The suppression factor relative to single crossing is
  $P_{\mathrm{single}}/P_{\mathrm{sustain}\text{-}K} = q^{-(K-1)}(1+o(1))$ as $q\to0$.
\item[\textnormal{(c)}] Let a true transition occur at checkpoint $t^{\ast}$ with
  post-transition per-checkpoint exceedance probability $q_1$ (operationally, the
  probability under the alternative that the de-noised readout exceeds $\tau$ at a
  post-transition checkpoint) over a window $W=M-t^{\ast}$. Then the sustained-$K$ miss
  probability (the false-negative rate, FNR) obeys the disjoint-block bound
  \[
    \mathrm{FNR}\leq (1-q_1^{K})^{\lfloor W/K\rfloor},
  \]
  and the expected detection latency is
  $\Delta\cdot(1-q_1^{K})/\bigl(q_1^{K}(1-q_1)\bigr)$ checkpoints, where $\Delta$ is the
  evaluation interval. Raising $K$ buys $q^{-(K-1)}$ false-positive suppression at the cost
  of latency growing like $q_1^{-K}$ and $\mathrm{FNR}$ rising when the true signal sits
  near $\tau$ (i.e.\ $q_1$ bounded away from $1$).
\end{enumerate}
\end{lemma}

Parts~(a) and~(b) are classical run statistics
\citep{feller1968,fulou2003runs,glaz2001scan}; part~(c), which maps the sustained criterion
onto a testbed's true-transition regime through the operationally defined $q_1$, is the new
general content of this reduction.

\begin{proof}[Proof of \textnormal{(c)}]
Partition the post-transition window into $\lfloor W/K\rfloor$ disjoint blocks of $K$
consecutive checkpoints. Under independence each block is entirely above $\tau$ with
probability $q_1^{K}$; the detector misses only if every block fails, giving
$\mathrm{FNR}\leq\prod(1-q_1^{K})=(1-q_1^{K})^{\lfloor W/K\rfloor}$ (an upper bound, since a
qualifying run may also straddle a block boundary). Detection waits for the first block of
$K$ consecutive exceedances; the number of Bernoulli$(q_1)$ trials to the first length-$K$
run has mean $(1-q_1^{K})/\bigl(q_1^{K}(1-q_1)\bigr)$, and multiplying by the interval
$\Delta$ gives the expected latency.
\end{proof}

\paragraph{The ``emergence rises with horizon'' pattern, demoted to a falsifiable prediction.}
The next statement isolates the one leg of the account the paper does not rest on.

\begin{corollary}[Nuisance-knob noise inflation]\label{cor:tdep}
Suppose a nuisance knob (for instance the in-context horizon $T$) leaves the null accuracy
$p$ unchanged but inflates the per-checkpoint noise scale $\sigma(T)$. Whenever $\tau>p$ and
the noise law is symmetric and log-concave (the Gaussian case being immediate),
$q(T)=\bar\Phi\bigl((\tau-p)/\sigma(T)\bigr)$ is increasing in $\sigma$, so
$P_{\mathrm{single}}(T)$ increases in $T$ with no change in the underlying computation,
while $P_{\mathrm{sustain}\text{-}K}$ (scaling as $q^{K}$) is far less sensitive.
\end{corollary}

The analytic core---$q$ increasing in $\sigma$ implies $P_{\mathrm{single}}$ increasing in
$\sigma$---is elementary. But the antecedent ``$\sigma$ grows with $T$'' is not
implied by the paper's own $T$-independent binomial noise $\sigma=\sqrt{p(1-p)/N}$; it
requires an auxiliary mechanism, either a decreasing effective sample size
$N_{\mathrm{eff}}(T)$ (longer sequences yielding fewer independent decision points or
correlated per-sequence outcomes) or a baseline $p(T)$ drifting toward $\tau$. The $T$-grid
logs adjudicate between the two. The noise-inflation route is inactive: the within-run
standard deviation is flat in $T$ (medians $0.091/0.090/0.088$ for $T=10/20/40$; Spearman
$r=-0.24$, $p=0.11$). The baseline-drift route is strongly active: the per-run mean accuracy
rises toward the threshold (medians $0.505/0.576/0.634$; Spearman $r=+0.94$, $p<0.001$), and
with it the single-crossing rate climbs from $1/15$ ($T{=}10$) through $10/15$ ($T{=}20$) to
$15/15$ ($T{=}40$) at the illustrative $\tau=0.8$ bar while the sustained-$K$ rate stays at
$0/45$. The ``emergence rises with
horizon'' pattern is therefore real in this testbed but runs through $p(T)\to\tau$ rather
than through the $\sigma(T)$ antecedent of Corollary~\ref{cor:tdep}, which remains an
unexercised prediction here; the single-versus-sustained asymmetry it forecasts is exactly
what the measured crossing rates show. The paper does not rest on this leg.

\paragraph{The pre-registered $\tau=0.7$ bar leads the empirical chain.}
The in-context-TD testbed of Section~\ref{sec:r2} operates two thresholds: its training
logs' online detector was pre-registered at $\tau=0.7$, and the calibration audit
additionally reports $\tau=0.8$, chosen so that a single case-study figure can show both
criteria operating (Section~\ref{sec:r2}). To keep the two tracks unambiguous:
$\tau=0.7$ is the pre-registered bar---fixed, together with the detection criteria, in
the analysis plan before the calibration analysis was run; we use ``pre-registered''
throughout in this analysis-plan sense, the plans being internal rather than lodged with
a third-party registry---and it carries the paper's headline numbers, whereas $\tau=0.8$
is an illustrative bar chosen for figure legibility, because the single-crossing and $K=2$
criteria are both visible at that bar. At the pre-registered bar the mechanism is at
its most severe, so we state that chain first, computed on the same 45 non-learning runs
($M=201$ checkpoints each). The measured pooled per-checkpoint exceedance is
$\bar q=0.099$---an order of magnitude above the $0.010$ measured at $\tau=0.8$---and the
single-crossing detector is saturated: it fires in $42/45$ runs. Because the 45 runs are
heterogeneous in baseline accuracy, the correct preflight tool this close to the noise
band is the per-run mixture prediction $\sum_i\bigl[1-e^{-(M-K+1)\hat q_i^{K}}\bigr]$,
with $\hat q_i$ the measured per-run exceedance rate. The mixture is estimated on the
same $45$ runs it predicts---an in-sample consistency check rather than an out-of-sample
forecast; its only out-of-sample support in this paper is the five-fold cross-validation
on the sixty-seed $\gamma=0.9$ grid (Section~\ref{sec:ladder}). It reproduces the entire observed
sustained-$K$ ladder to within $1.5$ of $45$ runs: observed (predicted) counts are $42$
($40.6$) for single crossings, $24$ ($24.5$) at $K{=}2$, $14$ ($13.6$) at $K{=}3$, $4$
($5.5$) at $K{=}4$, and $1$ ($1.5$) at $K{=}5$. The pooled-$\bar q$ nomogram reading of
Section~\ref{sec:nomogram} ($K\geq3.6$, i.e.\ $K=4$) under-prescribes once this
heterogeneity is priced in---the mixture still predicts $5.5/45\approx0.12$ false
positives per run at $K{=}4$---so the mixture is the prescribing tool at such a bar, and
it prescribes $K=5$, the smallest $K$ whose predicted per-run rate ($1.5/45\approx0.034$)
meets the $\alpha=0.05$ budget. At $K=5$ exactly one of 45 runs survives, against $1.5$
predicted: the pre-registered bar turns a $42/45$ false-positive saturation into a
calibrated $1/45$, an observed ${\sim}42\times$ reduction (mixture-predicted
${\sim}27\times$). Two lessons carry through the paper. First, noise suppression is
threshold-dependent: the required run length grows from $K=2$ at $\tau=0.8$ to $K=5$ at
$\tau=0.7$. Second, a practitioner applying Box~1 at a bar close to the noise band must
use the per-run mixture rather than a pooled-$q$ reading. The single-binomial reading is
worse still than the pooled one at this bar: it prescribes $K=3$, whose observed
false-positive rate is $14/45$. That failure is worked through as a cautionary example
after Box~1 (p.~\pageref{par:box1caution}), and it is why the checklist carries an
explicit mixture step. The $\tau=0.8$ computation
below is then the clean closed-form illustration of the same mechanism, at a bar where
the exact binomial tail and the i.i.d.\ formulas can be followed digit by digit.

\begin{corollary}[Recovering the in-context-TD numbers]\label{cor:numbers}
For the in-context-TD testbed of Section~\ref{sec:r2} ($N=32$ held-out sequences, $M=201$
checkpoints at steps $0,50,\dots,10{,}000$, chance baseline $p=0.562$, threshold
$\tau=0.8$), the exceedance probability is
$q(0.8)=P\!\bigl(\mathrm{Binom}(32,0.562)\geq 26\bigr)=0.00272$. Theorem~\ref{thm:single}
then gives $P_{\mathrm{single}}=1-(1-0.00272)^{201}=0.421$ (expected $18.9/45$ runs), and
Lemma~\ref{lem:sustain} gives $P_{\mathrm{sustain}\text{-}2}=200\cdot(0.00272)^{2}=0.00147$
(expected $0.066\approx0.07/45$, matching the exact finite-Markov-chain value), a
suppression factor of $0.421/0.00147\approx 286\times$.
\end{corollary}

\paragraph{Empirical validation of the noise model.}
With $p=0.562$ and $N=32$ the predicted per-checkpoint standard error is
$\sigma=\sqrt{0.562\times0.438/32}=0.0877$; the measured within-run standard deviation of
validation accuracy is $0.089$ (median over the 45 per-run logs), matching to within $2\%$.
This validates the marginal Bernoulli/Gaussian readout empirically---a necessary, not
sufficient, condition (it does not validate independence). At the pre-registered threshold
$\tau=0.7$, $\lceil0.7\times32\rceil=23$ and $q(0.7)=0.0517$ (exact binomial at the median
baseline), giving $P_{\mathrm{single}}(0.7)\approx1.000$ (expected $45/45$)---consistent
with the observed $42/45$ saturation reported above. (The measured pooled
$\bar q=0.099$ exceeds this median-baseline value because $q$ is convex in the baseline
accuracy; that heterogeneity is exactly what the per-run mixture prices.)

\paragraph{Honesty caveats on the illustrative-bar suppression factor and the $26/45$.}
The suppression factor of Corollary~\ref{cor:numbers} and the $0.07/45$ sustained prediction hold under the
positive-association hypothesis of Theorem~\ref{thm:single}(b) with negligible
autocorrelation. The autocorrelation is now measured (Remark~\ref{rem:assoc}): the median
detrended lag-1 value is $+0.03$, with $41/45$ runs inside the white-noise band. Even
pricing that residual dependence conservatively---replacing one factor of $q$ in
$(M-1)q^{2}$ by the linearly inflated conditional $q+\rho_{1}(1-q)$ at the median
$\rho_{1}$---the suppression factor remains ${\approx}24\times$, and the observed sustained-2
count is $0/45$, consistent with the i.i.d.\ prediction of $0.066/45$. The factor is
thus the i.i.d.\ value with measured-small dependence, bounded below by ${\sim}24\times$
under a deliberately pessimistic dependence model.
That floor is pessimistic in a way the data can quantify. The quantity the run-length
formulas actually depend on is the autocorrelation of the exceedance indicator
$\mathbf{1}\{a_{m}\geq\tau\}$, not of the accuracy trace, and it is smaller: measured at
lags $1$ through $5$ across the 45 runs, the mean indicator autocorrelation is $+0.0002$
with every lag satisfying $|\rho_{\ell}|<0.03$. The i.i.d.\ assumption behind
Theorem~\ref{thm:single} is therefore better supported than the ${\approx}24\times$ floor
implies; we report the floor as a conservative bound rather than a best estimate.
The observed single crossings are $26/45$, above the $18.9/45$ predicted at the median $p$:
the 45 cells are heterogeneous in baseline accuracy (finals reach $0.875$) and $q$ is convex
in $p$. The direct mixture computation over the empirical per-run baselines,
$\mathbb{E}_{\hat p}\bigl[P_{\mathrm{single}}(\hat p)\bigr]$ with $\hat p$ the per-run mean
accuracy, gives $0.541$ (expected $24.3/45$), and the observed $26/45$ lies well within one
binomial standard deviation ($\pm3.3$) of that prediction---so the excess over $18.9/45$ is
quantitatively accounted for by baseline heterogeneity, not by a mechanism outside the noise
model. All $q$ values in this section are exact binomial tail probabilities
($q(0.8)=0.002716$, $q(0.7)=0.0517$); the normal-$z$ figures quoted for intuition agree to
the printed digits.


\section{Calibration nomogram for sustained-crossing criteria}\label{sec:nomogram}

The sustained-$K$ criterion becomes a reusable instrument once $K$ is chosen to meet a
target false-positive budget. From the small-$q$ form of Lemma~\ref{lem:sustain}(a), a
target noise-only budget $\alpha$ requires $(M-K+1)q^{K}\leq\alpha$, i.e.
\begin{equation}\label{eq:nomogram}
  K \;\geq\; \frac{\ln M - \ln\alpha}{\ln(1/q)}.
\end{equation}
For the paper's parameters ($\alpha=0.05$, $M=201$, $q=0.00272$),
\[
  K \geq \frac{\ln 201 - \ln 0.05}{\ln(1/0.00272)}
      = \frac{5.30 + 3.00}{5.91}
      = 1.40,
\]
so $K=2$ suffices ($P_{\mathrm{sustain}\text{-}2}=0.00147\ll0.05$). The $K=2$ criterion used
throughout Section~\ref{sec:r2} is thus derived from Eq.~\eqref{eq:nomogram} rather
than assumed. Figure~\ref{fig:nomogram} renders the full instrument: contours of the
required $K$ over the $(q,\alpha,M)$ ranges a practitioner encounters, with this paper's
operating point marked.

\begin{figure}[htbp]
\centering
\figtikz[\linewidth]{E2_nomogram}
\caption{$K$-selection nomogram. Required sustained-run length $K$ from
Eq.~\eqref{eq:nomogram} as a function of the per-checkpoint exceedance probability $q$
(both axes logarithmic), with contours for target noise-only false-positive budgets
$\alpha\in\{0.05,0.01,0.001\}$ (colour) at $M\in\{50,200,1000\}$ checkpoints (line type).
A practitioner computes $q$ from $(N,p,\tau)$ and reads off the minimum $K$ meeting the
budget $\alpha$ over $M$ checkpoints. The star marks the in-context-TD operating point
($q=0.00272$, $M=201$), which reads $K\geq2$ at $\alpha=0.05$; the dotted horizontal line
marks $K=2$.
(b)~Observed versus null-predicted single-crossing fires on three corpora: the
instrument predicts the false-positive count before any correction is applied
(in-context TD $42$ observed; Pythia $4$ observed against $5.7$ binomial-null
predicted; BIG-Bench $40$ observed against $17.7$ expected).
(c)~The pooled-$q$ cautionary example: the naive single-pooled route prescribes
$K=3$ and observes $14/45$ false positives ($6.2\times$ over the $0.05$
budget), while the per-run-mixture route prescribes $K=5$ and restores
calibration at $1/45$.
(d)~Suppression-versus-latency tradeoff: the exact $286\times$ suppression at
the operating point ($q=0.00272$, $K=2$) costs one checkpoint of confirmation
latency; on the genuine-emergence controls sustained-$K$ removes every false
fire while keeping $8/9$ genuine detections at zero added latency.
Panel summaries: (a)~K nomogram. (b)~calibration check. (c)~pooled-q caution. (d)~power-latency tradeoff.}
\label{fig:nomogram}
\end{figure}

Box~1 condenses the correction into an operational calibration checklist.

\begin{tcolorbox}[colback=gray!6, colframe=gray!40, arc=2pt,
  left=6pt, right=6pt, top=4pt, bottom=4pt,
  title={\small\bfseries Box 1. Calibration checklist: before declaring emergence from a threshold crossing},
  fonttitle=\small\bfseries, breakable, enhanced]
\begin{enumerate}
\item \textbf{Compute $q$} $=P(\mathrm{Binom}(N,p)\geq\lceil\tau N\rceil)$ from your actual
  $N$, the baseline $p$ your trace plateaus at, and threshold $\tau$. Use the plateau, not
  the task's chance level: if the model learns an easier component of the task, the two
  differ and only the plateau predicts crossings. If $qM\gtrsim0.1$, single crossings are
  noise.
\item \textbf{Validate the noise model:} does the empirical within-run $\sigma$ match
  $\sqrt{p(1-p)/N}$? (Necessary, not sufficient.)
\item \textbf{Check autocorrelation $\rho_\ell$ out to $\ell=K$.} If $\rho_\ell>0$, your
  sustained-$K$ false-positive rate is worse than the i.i.d.\ formula.
\item \textbf{If per-run baseline levels vary, price the heterogeneity.} The trigger:
  replace the single binomial $q$ of step~1 with the per-run mixture
  $\sum_i\bigl[1-\exp(-(M-K+1)\hat q_i^{\,K})\bigr]$ over the measured per-run exceedance
  rates $\hat q_i$ whenever $\bar qM \gtrsim 1$ or the per-run rates span more than
  about a factor of two of their mean. This mixture is the quantity the nomogram
  actually prices; routing through the single binomial $q$ when $\bar qM$ is not
  small under-prescribes $K$, because $q$ is convex in the baseline accuracy and the
  runs above the pooled mean dominate the false-positive count. (The factor-of-two
  span and the $\bar qM \lesssim 1$ sufficiency zone are empirical, observed in the
  testbeds of Section~3; a rigorous bound is open.)
\item \textbf{Pick $K$} from Eq.~\eqref{eq:nomogram} for your target $\alpha$, evaluated on
  whichever of step~1 or step~4 applies.
\item \textbf{Freeze nuisance knobs} (horizon $T$, optimizer): if a ``trend'' tracks a knob
  that inflates $\sigma$, it may be a noise artifact---verify that the sustained rate is
  flat while the single rate rises.
\item \textbf{Include a solved-reference cell} so that a real transition would actually
  cross.
\end{enumerate}
\end{tcolorbox}

\paragraph{A worked cautionary example: the checklist applied naively to this paper's own
testbed.}\label{par:box1caution}
Step~4 is not a refinement, and the cost of skipping it is measurable on the testbed of
Section~\ref{sec:r2}. At the pre-registered $\tau=0.7$ bar, step~1 evaluated at the median
baseline gives $q=0.0517$, and Eq.~\eqref{eq:nomogram} with $M=201$ and $\alpha=0.05$ then
reads $K\geq2.80$, i.e.\ $K=3$. The observed false-positive rate of that prescription is
$14/45=0.311$, which is $6.2$ times the nominal budget. The per-run mixture of step~4,
computed on the same runs, prescribes $K=5$ instead and restores calibration: $1/45$
observed against $1.5$ predicted, inside $\alpha$. The pooled-$\bar q$ reading
($\bar q=0.099$) lands in between at $K=4$, still over budget at $4/45=0.089$. Two
properties of this bar drive the gap. The baseline heterogeneity across the 45 runs is
large (per-run $\hat q_i$ span $0.004$ to $0.133$), and $\bar qM\approx20$ is far from the
small-$\bar qM$ regime in which a single $q$ summarises the ensemble. The naive reading is
the failure mode this checklist exists to catch, and reporting it here fixes the operating
range over which a single binomial $q$ can be trusted: $\bar qM\lesssim1$, with per-run
rates within a factor of about two of their mean.

\section{Methods}\label{sec:methods}

The three case studies come from a shared suite of tiny-transformer emergence experiments,
all with analysis plans (detector thresholds, criteria, and outcome branches) fixed before
analysis; for the re-analysis corpus of Section~\ref{sec:r3} the runs pre-existed and only
the analysis plan was fixed in advance. We summarise the common design here; the
testbed-specific sweeps (ordering policies, horizons, context lengths, seed counts) are
stated in each case study below.

\paragraph{Testbeds.} (1)~An in-context copy/induction family---the one task in the suite
that is learnable, non-Fourier, and shows a measurable emergence delay---used to test whether
data-presentation order compresses the emergence delay (Section~\ref{sec:r1}). (2)~An
in-context temporal-difference (TD) value-learning testbed, used here as the primary
validation of Section~\ref{sec:theorem} rather than a solved trainable-control result
(Section~\ref{sec:r2}). (3)~A post-hoc re-analysis of 364 existing runs across 11 closed
experiment namespaces, used to test early-trajectory outcome prediction
(Section~\ref{sec:r3}).

\paragraph{Emergence readout.} Emergence is read by a threshold-crossing detector (first
crossing of a fixed score threshold) on a fixed-step evaluation grid, matching the
operationalizations cited in Section~\ref{sec:intro}. Because a single-crossing
rule can fire on transient noise, we contrast it throughout against a sustained-emergence
criterion and a solved-reference (oracle) cell.

\paragraph{Optimizers and seeds.} Experiments sweep the \{AdamW, SGD with momentum (SGDM),
Muon (an orthogonalized-momentum optimizer)\} optimizers with multiple seeds per cell; the
exact grid for each testbed is given in the corresponding case study.

\section{Validation case studies}\label{sec:results}

The reduction of Section~\ref{sec:theorem} makes three checkable predictions on a shared
tiny-transformer programme. The primary validation is the in-context-TD grid
(Section~\ref{sec:r2}), which instantiates Corollary~\ref{cor:numbers} to within $2\%$; two
further case studies (Sections~\ref{sec:r1} and~\ref{sec:r3}) exercise the
floor-censoring/learnable-target and nuisance-axis disciplines the reduction implies. In
each case the null is the theorem's empirical anchor rather than the contribution, which is
why its individual weakness (a borderline $p$, a non-learning grid, a power-bounded
comparison) is acceptable. Two external-data case studies then apply the same instrument
beyond the programme: the public BIG-Bench release (Section~\ref{sec:bigbench}) and real
Pythia pretraining checkpoints (Section~\ref{sec:pythia}).

\subsection{Primary validation: in-context TD does not learn, and single crossings fire on noise}\label{sec:r2}

This grid is the primary validation of Section~\ref{sec:theorem}: it instantiates
Corollary~\ref{cor:numbers} to within $2\%$ and exhibits the single-versus-sustained gap the
reduction predicts.
On an in-context temporal-difference value-learning testbed that remains a calibration vignette rather than a solved trainable-control result (3 optimizers $\times$ horizon
$T \in \{10, 20, 40\}$ $\times$ 5 seeds), the model does not learn the task: final
validation accuracy plateaus far below competence almost everywhere (median $0.562$, min $0.406$, max
$0.875$; only 1/45 cells exceed $0.8$), and per-evaluation trajectories oscillate
within a typical band ($0.5 \leftrightarrow 0.72$) rather than showing a sustained jump
(worst-case single-cell excursions reach wider, e.g.\ $0.344$--$0.844$ in the case study below).
At the pre-registered online bar $\tau=0.7$ the single-crossing detector fires in $42/45$
of these non-learning runs, and the mixture-prescribed sustained-$5$ criterion reduces
this to $1/45$ (Section~\ref{sec:theorem}); the audit below is reported at $\tau=0.8$,
where the prescribed correction is $K=2$, so that a single figure can show both criteria
operating. At that bar, a threshold-crossing
emergence detector (first crossing of $\tau=0.8$) likewise fires on transient noise crossings:
cells flagged ``emerged'' settle back to $\approx 0.56$. The apparent pattern---
``emergence rate rises with horizon (T10: 0--1/5 $\to$ T40: 5/5) and emergence step drops
with horizon''---is an artifact of the noise detector (larger $T$ makes the eval estimate
noisier and the spurious $0.8$-crossing earlier), the falsifiable prediction of
Corollary~\ref{cor:tdep}, not evidence of faster TD computation.
None of the headline questions (sharp emergence, precursor diagnostics,
optimizer-moves-the-transition) is adjudicable until the testbed is recalibrated. The
recipe: add a solved-reference cell (here, an exact Bellman-oracle reference under the
same generator, which reaches strict accuracy $1.000$), use a de-noised readout
(smoothed / held-out accuracy), and require a sustained-emergence criterion (accuracy held
above threshold for a window), not a single crossing. As a calibration boundary, the same
transformer can fit finite Bellman-labelled support sets above $0.8$ in 5/6 cells while
fresh Markov reward process (MRP) generalisation accuracy remains far below the
positive-control bar (maximum $0.605$ vs the $0.8$ bar), establishing
that fitting labelled support is not TD emergence (Fig.~\ref{fig:supervised-fit}).

\begin{figure}[htbp]
\centering
\figtikz[\linewidth]{E2_supervised_fit}
\caption{Supervised Bellman-label fit sanity check. Finite-support Bellman labels often
exceed 0.8 while fresh-MRP accuracy stays low, so fitting labelled support is not TD emergence.
The same transformer and evaluator can fit finite Bellman-labelled support sets above 0.8
in 5/6 cells (maximum 0.898), but fresh-MRP (generalisation) accuracy remains far below
the sustained positive-control bar (maximum 0.605), establishing a calibration boundary
between memorisation of labelled support and genuine in-context TD emergence.
Panel summaries: (a)~finite-support fit at support $128$. (b)~finite-support fit at support $512$. (c)~support-fit margin over fresh MRP. (d)~support-fit threshold tally.}
\label{fig:supervised-fit}
\end{figure}

The non-signal is itself
informative: Figure~\ref{fig:td-grid} shows the full grid of per-evaluation accuracy
trajectories, which oscillate around the plateau with transient spikes that fall back, alongside
the precursor metric (the Bellman residual) that never descends toward zero---visual
confirmation that the threshold detector is firing on noise rather than on a forming value
function.
Table~\ref{tab:tdcalib} reports the
calibration audit and Table~\ref{tab:tdgrid} gives the per-cell quantitative breakdown:
transient $0.8$ crossings occur in 26/45 original runs, but none survive
a 2-evaluation or 5-evaluation sustained criterion. A second de-noised T40 rerun with
larger held-out evaluation batches also produces no $0.8$ crossing at all, strengthening the
negative calibration reading.

\paragraph{The grid's trivial floor, and where its null parameter comes from.}
Calling this grid chance-level, as an earlier reading of it did, overstates the negative.
The task's trivial floor is the accuracy of a constant predictor that always emits the
modal value bucket, and that floor can be computed exactly by simulating the generator.
For the grid's configuration (ten states, $\gamma=0.9$, nine value buckets over
$[-10,10]$; $4\times10^{5}$ sampled Markov reward processes) the target-bucket marginal is
$0.000/0.003/0.056/0.247/0.390/0.246/0.056/0.003/0.000$, so the constant predictor scores
$0.390$. The grid's median final accuracy of $0.562$ sits $0.17$ above that floor. The
model therefore does extract real structure---the same context-mean regression that the
$\gamma=0.9$ ladder rung of Section~\ref{sec:ladder} exhibits more strongly---and then
plateaus, without approaching either the task ceiling or the detection bar. The grid is a
partially learned, drifted baseline rather than a pure-noise trace, and we label it that
way throughout. The substantive reading is unchanged: the model never acquires
temporal-difference computation, so the crossings remain false positives for emergence,
and the Bellman-residual precursor of Fig.~\ref{fig:td-grid} never descends.

The null parameter $p=0.562$ of Corollary~\ref{cor:numbers} is accordingly not the task's
chance level but an empirical per-checkpoint exceedance estimate taken from the runs being
tested. That is the quantity the scan calculation needs. Theorem~\ref{thm:single} is a
statement about how often a stationary trace crosses a bar, so its parameter must be the
level the trace actually sits at; substituting the trivial floor $0.390$ would give
$q(0.8)=1.3\times10^{-6}$ and $0.014$ expected single crossings in $45$ runs against the
$26$ observed, a prediction calibrated against a model that is not the one being run.
Estimating $p$ from the runs makes the false-positive computation self-consistent, at the
cost of making it descriptive rather than a priori. A practitioner applying Box~1 before
training has no such estimate and must instead bound $p$ by the highest plateau the model
could plausibly reach, which is one of the things the solved-reference cell of step~7 is
for.

\paragraph{Case study: a single false-positive emergence.}

Figure~\ref{fig:case} shows the validation-accuracy trajectory of one representative cell
(AdamW, $T=40$, seed~3).
At step~150 the accuracy spikes to $0.844$, crossing a single-crossing detector threshold;
the detector therefore flags this cell as ``emerged.'' (The calibration audit reported throughout
this case study fires at $\tau=0.8$; the original in-context-TD training logs additionally record an
online emergence field at the pre-registered $\tau=0.7$. At that lower bar the exceedance
rate is tenfold higher and $K=2$ no longer suffices---the full $\tau=0.7$ sustained-$K$
ladder, with the mixture-prescribed $K=5$, leads the empirical chain of
Section~\ref{sec:theorem}. The case study uses $\tau=0.8$, where $K=2$ is the prescribed
correction, so that a single figure can show both criteria operating.)
Over the remaining $9{,}850$ steps the
accuracy oscillates between $0.344$ and $0.844$ without ever stabilising above threshold,
and the run ends at $0.5625$---within rounding of the grid's plateau level of $0.562$
(the median across all 45 cells). This is a textbook false positive: a single noisy
upward fluctuation triggers the detector, while the model never acquires the task. The
recalibration recipe above (a solved-reference cell, a de-noised readout, and a
sustained-emergence criterion) would prevent all three failure modes simultaneously: a
solved-reference cell establishes that a sustained crossing is achievable; de-noising
damps the spike that causes the first crossing; and requiring the threshold to be held for
a window rules out the isolated crossing at step~150.

\begin{figure}[htbp]
\centering
\figtikz[0.97\linewidth]{E2_td_grid}
\caption{The in-context-TD calibration grid shows no emergence to detect.
(a)~Per-evaluation validation-accuracy trajectories for all nine optimizer $\times$ horizon
cells (five seeds each, 45 runs); curves oscillate around the grid's plateau level (dotted) with
transient excursions that fall back rather than a sustained jump. The dashed line marks the
$0.8$ detector threshold and the lower dotted-green line the $0.7$ threshold; transient
single crossings become more frequent at larger horizon $T$ because the per-evaluation
estimate is noisier, not because computation is faster.
(b)~Single-crossing fire rate at $\tau=0.7$ per cell ($42/45$ runs): the
false-positive count the run statistics predict; bar centres sit close enough that
per-bar labels sit inside the bars rather than above them.
(c)~The precursor metric (median Bellman residual over training, per optimizer) never
descends toward zero, confirming that no value function is being learned for a detector to
fire on.
(d)~The sustained-$K$ ladder: observed false-positive counts collapse with $K$
($42/24/14/1$ at $\tau=0.7$; $26/0/0/0$ at $\tau=0.8$), with the per-run
mixture predicting $1.5$ at $K=5$ against the observed $1$.
Panel summaries: (a)~validation grid. (b)~single crossings. (c)~precursor null. (d)~mixture calibration.}
\label{fig:td-grid}
\end{figure}

\begin{figure}[htbp]
\centering
\figtikz[\linewidth]{E2_case}
\caption{A single false-positive emergence makes the false-positive account of
Section~\ref{sec:theorem} (Corollary~\ref{cor:numbers})
visible on one trajectory (AdamW, $T=40$, seed~3). The validation accuracy spikes to a
peak of $0.844$ at step~150 (orange diamond), crossing the $\tau=0.8$ detector threshold
(dashed red); over the full run it crosses the threshold at six isolated points (open red
circles), so the single-crossing detector fires. None of these crossings is followed
by a second consecutive evaluation above threshold, so the sustained-$K{=}2$
detector records zero qualifying pairs and rejects the cell. The run oscillates within the
plateau band ($0.5625 \pm 1$ binomial SE, grey) and settles by step~10\,000 to $0.5625$,
the grid's plateau level and $0.17$ above the task's constant-predictor floor of $0.390$. This is the per-run picture behind the population suppression factor of
${\approx}286\times$ (under the association hypothesis of Theorem~\ref{thm:single}): a
single noisy upward fluctuation triggers the single-crossing
detector, whereas requiring the threshold to be held for one further evaluation removes the
false positive.
(b)~Zoom on the crossing cluster: the firing is a transient spike, not a
regime change.
(c)~Longest run at the bar versus the geometric null at the measured
per-checkpoint exceedance rate: the observed run length is unremarkable under
the memoryless null.
(d)~Sustained-$K$ rejection: the same trajectory fires the single-crossing
detector and is silent at $K\geq2$.
Panel summaries: (a)~false-positive trajectory. (b)~exceedance clustering. (c)~geometric null. (d)~sustained-K rejection.}
\label{fig:case}
\end{figure}

\begin{table}[htbp]
\centering
\caption{Calibration audit for the in-context-TD grid, at both operating thresholds:
the pre-registered online bar $\tau=0.7$ (mixture-prescribed $K=5$) and the illustrative
audit bar $\tau=0.8$ ($K=2$). The exact Bellman oracle is a
solved-reference readout for the same data generator/evaluation target; the original trained
grid is the 45-run optimizer $\times$ horizon sweep, and the de-noised rerun is a bounded T40
check with larger held-out evaluation batches.}
\label{tab:tdcalib}
\begin{threeparttable}
\begin{tabular}{llcl}
\toprule
Readout & Criterion & Count/value & Interpretation \\
\midrule
Exact Bellman oracle\tnote{a} & strict acc. & 1.000 & solved reference \\
Trained grid & median final acc. & 0.562 & 45 runs \\
Single crossing & $\geq0.7$ once & 42/45 & pre-registered bar; saturated \\
Sustained crossing & $\geq0.7$ for 5 evals & 1/45 & mixture-prescribed $K{=}5$ \\
Single crossing & $\geq0.8$ once & 26/45 & spike-prone \\
Sustained crossing & $\geq0.8$ for 2 evals & 0/45 & de-noised \\
Sustained crossing & $\geq0.8$ for 5 evals & 0/45 & primary guard \\
De-noised T40 rerun\tnote{b} & $\geq0.8$ once / smooth5 & 0/9 / 0/9 & no rescue \\
\bottomrule
\end{tabular}
\begin{tablenotes}[flushleft]
\footnotesize
\item[a] Analytic exact solution: the evaluator queries the true Bellman value $V^\pi$ for each MRP directly via dynamic programming; this is not a trained run.
\item[b] 9 held-out runs (3 optimizers $\times$ 3 seeds); re-run with larger held-out evaluation batches (128 sequences); no $\geq0.8$ crossing under the de-noised readout.
\end{tablenotes}
\end{threeparttable}
\end{table}

\begin{table}[htbp]
\centering
\caption{In-context TD calibration grid: median final validation accuracy and threshold-crossing counts
per optimizer $\times$ horizon ($n=5$ seeds each; 45 runs total).
Single crossing: at least one evaluation point $\geq0.8$.
Sustained-2: $\geq0.8$ held for $\geq2$ consecutive evaluations.
No cell achieves a sustained crossing; the single-crossing pattern grows with $T$ (detector artifact, not faster computation).}
\label{tab:tdgrid}
\small
\begin{tabular}{llccc}
\toprule
Optimizer & Horizon $T$ & Median final acc. & Single cross.\ $\geq0.8$ & Sustained-2 $\geq0.8$ \\
\midrule
AdamW & 10 & 0.594 & 0/5 & 0/5 \\
AdamW & 20 & 0.531 & 2/5 & 0/5 \\
AdamW & 40 & 0.594 & 5/5 & 0/5 \\
\addlinespace
Muon  & 10 & 0.500 & 1/5 & 0/5 \\
Muon  & 20 & 0.594 & 3/5 & 0/5 \\
Muon  & 40 & 0.563 & 5/5 & 0/5 \\
\addlinespace
SGDM  & 10 & 0.594 & 0/5 & 0/5 \\
SGDM  & 20 & 0.563 & 5/5 & 0/5 \\
SGDM  & 40 & 0.594 & 5/5 & 0/5 \\
\midrule
\multicolumn{2}{l}{Aggregate (45 runs)} & 0.562 & 26/45 & 0/45 \\
\bottomrule
\end{tabular}
\end{table}

\subsection{Discount ladder: no learnable rung, with validation under baseline
drift}\label{sec:ladder}

A natural objection to the grid above is selection: perhaps the testbed was simply an
unlearnable cell, and a slightly easier variant would have emerged. We therefore built a
difficulty ladder inside the same pipeline---identical model, tokenization, evaluator,
and detector; six states, horizon $T=40$; only the discount changes,
$\gamma\in\{0,0.3,0.5,0.7,0.9\}$ (five seeds at $\gamma\in\{0,0.3\}$, ten at
$\gamma\in\{0.5,0.7\}$, sixty at $\gamma=0.9$). At $\gamma=0$ the value function collapses
to the immediate discretized reward, $V(s)=r(s)$, so the task degenerates to an
in-context lookup of the queried state's reward token---the same competence class as the
induction positive control above. The information-theoretic ceiling is $0.95$--$0.96$ at
every rung (an oracle restricted to in-context evidence), and the guess-the-modal-bucket
baselines are $0.149/0.159/0.305$ for $\gamma=0/0.3/0.9$.

Two results. First, no rung is learnable within the $10^4$-step budget: from $\gamma=0$
through $\gamma=0.7$ the best checkpoint accuracy rises only gently with $\gamma$ (medians
$0.500$, $0.531$, $0.562$, $0.656$ across the four rungs)---above the modal baselines, so
some signal is extracted, but far below both the ceiling and the detection bar; neither
detector ever fires ($0/30$ runs across the four rungs). Even the lookup rung does not
train at this budget, which localizes the family's bottleneck to state-resolved
in-context retrieval rather than to value bootstrapping, and rules out the
easier-variant objection at the studied scale: there is no solvable rung between the
canonical task and its induction-grade reduction that the grid's budget could have
reached. (Quadrupling the budget shows the lookup rung is learnable but very late; we
return to this in Section~\ref{sec:mnaxes}, where the late transition serves as an
in-family solved reference.)

Second, the $\gamma=0.9$ rung turns baseline drift from a theoretical mechanism into a
live measurement---and, at sixty seeds, into a rate-level calibration test. There the
achievable component is context-mean regression (values concentrate, hence the $0.305$
modal baseline), and the model genuinely learns it, lifting the running accuracy
baseline to $\bar p = 0.642$ without ever reaching sustained competence (best checkpoint
$0.813$--$0.938$, bar $0.8$, ceiling $0.96$). The detectors respond exactly as
Section~\ref{sec:theorem} predicts. Single-crossing fires in $60/60$ seeds. The per-run
mixture prediction (measured per-run exceedances $\hat q_i$, pooled $\bar q=0.034$)
gives $13.2$ expected sustained-2 firings in 60 runs; $15$ are observed (rate $0.250$,
$95\%$ CI $[0.147,0.379]$). The estimate is not circular: in a five-fold cross-validation
(per-run $\hat q$ estimated on four fifths of the seeds, the held-out fifth predicted),
the held-out counts run from $1$ to $5$ sustained-2 firings per fold against per-fold
predictions of $2.5$--$2.8$, depending on how the sixty seeds are partitioned
($3,2,3,3,4$ under a strided assignment on seed order; $3,1,5,3,3$ under a contiguous
one). Every fold lies within two standard deviations of its prediction under either
partition, and the fold totals agree with the pooled figures in all cases ($15$ observed
vs $13.2$ predicted). The nomogram of Section~\ref{sec:nomogram} prescribes
the repair before looking at the data: at $q=0.034$ it reads $K\geq2.45$, i.e.\
$K=3$---and sustained-3 fires once in $60$ seeds against a predicted count of $0.64$
(rate CI $[0.0004,0.089]$), within its $\alpha=0.05$ per-run budget. The criterion is not a universal constant but an instrument: when
partial learning (or any nuisance) drifts the baseline toward the bar, $K$ must be
re-read from the nomogram, and the re-read value restores calibration.

Third, the horizon effect of Section~\ref{sec:theorem} replicates on eighty runs that
share nothing with the main grid except the detector. The $\gamma=0.9$ rung was also run
at two shorter horizons, ten seeds each at $T=10$ and $T=20$, alongside the sixty-seed
$T=40$ arm; all use six states and AdamW, where the main grid used ten states and three
optimizers. The two grids agree closely on the quantity the account turns on. Median
per-run accuracy rises $0.509\to0.582\to0.642$ with $T$ here against
$0.505\to0.576\to0.634$ there, and the single-crossing count rises
$1/10\to9/10\to60/60$ against $1/15\to10/15\to15/15$ (Table~\ref{tab:horizonrep}). Since
the replication changes the state count, the optimizer set and the seeds, the horizon
effect is a property of how the baseline depends on $T$ rather than of one grid's
configuration.

The two grids differ in one cell, and the difference is the predicted one. At $T=40$ and
$\tau=0.8$ the replication records $15/60$ sustained-2 firings where the main grid records
$0/15$. The replication's baseline sits slightly higher ($\bar q=0.034$ against $0.025$),
which raises the i.i.d.\ expectation from $1.9$ per $15$ runs to $13.2$ per $60$; the
replication lands on its expectation while the main grid falls below it, the direction
Theorem~\ref{thm:single}(b) permits under positive association. Sustained-3 fires $0/10$,
$0/10$ and $1/60$ across the three arms. At the pre-registered $\tau=0.7$ bar the same
ordering appears earlier and more sharply, with single crossings at
$9/10\to10/10\to60/60$ and sustained-2 at $1/10\to7/10\to60/60$: once the baseline has
drifted this far, a sustained criterion carried over unchanged from a shorter horizon
saturates too, which is the case for re-reading $K$ from the nomogram at each operating
point rather than fixing it once.

\begin{table}[htbp]
\centering
\caption{Independent replication of the horizon effect on the $\gamma=0.9$ ladder rung
(six states, AdamW, $N=32$, $M=201$), to be compared with the main grid's
$1/15\to10/15\to15/15$ single-crossing ladder at $\tau=0.8$ (ten states, three
optimizers). Accuracy is the median across runs of each run's mean over all $201$
checkpoints. Detector counts are given at both operating bars.}
\label{tab:horizonrep}
\small
\begin{tabular}{lccccccc}
\toprule
& & & \multicolumn{2}{c}{$\tau=0.8$ (illustrative)} & \multicolumn{2}{c}{$\tau=0.7$ (pre-registered)} \\
\cmidrule(lr){4-5}\cmidrule(lr){6-7}
Horizon $T$ & Runs & Mean acc. & Single & Sustained-2 & Single & Sustained-2 \\
\midrule
10 & 10 & 0.509 & 1/10  & 0/10  & 9/10  & 1/10 \\
20 & 10 & 0.582 & 9/10  & 0/10  & 10/10 & 7/10 \\
40 & 60 & 0.642 & 60/60 & 15/60 & 60/60 & 60/60 \\
\bottomrule
\end{tabular}
\end{table}

Together with the induction-side validation above, the sustained criterion is now checked
in three settings, which span two drift levels rather than three separate regimes: a
mildly drifted baseline (the 45-run grid, $0.17$ above its constant-predictor floor,
$0/45$ at the illustrative $\tau=0.8$ bar), a strongly drifted one (this rung, $0.35$
above its floor, where a stale $K$ fails at the predicted rate and the
nomogram-corrected $K$ holds), and genuine transitions ($90/90$ detected at zero
latency). No arm of the TD family sits at its trivial floor. What makes the
non-transition arms noise-like for the detector is that their traces are stationary, not
that they are at chance, which is the sense in which Theorem~\ref{thm:single} applies to
them.

\subsection{Stress tests along the checkpoint and eval-size axes}\label{sec:mnaxes}

The false-positive formula has two free axes besides the threshold: the number of
checkpoints $M$ and the evaluation-set size $N$ (which sets $q$ through the binomial
tail). We measured both in the drifted regime ($\gamma=0.9$), where the predictions are
non-trivial.

Along the $N$ axis (five seeds each at $N=16$ and $64$, sixty at the default $N=32$;
$M=201$), the measured exceedance tracks the binomial-tail computation of Box~1 step~1
within a factor of two despite the drifting baseline: observed $\bar q = 0.133$, $0.034$,
$0.004$ against predicted $0.119$, $0.030$, $0.002$ at $N=16$, $32$, $64$ ($\bar p
\approx 0.64$). The detector consequences span the full range of failure and repair. At
$N=16$, single-crossing and sustained-2 fire in $5/5$ seeds and even sustained-3 fires
in $2/5$---which the nomogram itself predicts, since at $q=0.133$ it reads $K\geq4.1$:
a $K$ calibrated for $N=32$ is stale at $N=16$. At $N=64$, single-crossing itself drops
to $3/5$ and sustained-2 is silent. Evaluation-set size, not any training-side knob, is
the dominant false-positive lever in this regime, and the preflight computation
anticipates its effect quantitatively.

Along the $M$ axis the same rung was run at two longer budgets under the identical
configuration, ten seeds each: $M=401$ (doubled training) and $M=801$ (quadrupled). We
report both, since the axis is only informative if its endpoints are shown. Single-crossing
is saturated throughout, $60/60$ at $M=201$ and $10/10$ at both longer budgets, as it must
be once $Mq\gg1$. The sustained-2 counts are $15/60$, $2/10$ and $5/10$, against i.i.d.\
envelopes of $13.2$, $3.7$ and $6.6$ runs computed from each arm's own measured exceedance
($\bar q=0.034$, $0.034$, $0.038$); every arm lies at or below its envelope, the direction
Theorem~\ref{thm:single}(b) permits under positive association. Sustained-3 gives $1/60$,
$0/10$ and $1/10$ against $0.6$, $0.2$ and $0.5$ expected. The false-positive rate
therefore rises with $M$ at fixed $\tau$ and fixed configuration---the sustained-2 rate
roughly doubles from $0.25$ to $0.50$ as the grid is refined fourfold---which is
Theorem~\ref{thm:single}(a)'s $M$-monotonicity measured directly rather than assumed. None
of this is learning: at $M=801$ the drifted rung's accuracy is as flat as at $M=201$ (mean
over the first fifth of checkpoints $0.638$, over the last fifth $0.657$; no run holds the
bar for more than three consecutive checkpoints, and final accuracies stay between $0.594$
and $0.781$). Quadrupling the budget buys this rung no competence, only more opportunities
to cross. The lower rungs serve
as negative controls along the same axis: across the non-learning runs
($\gamma\in\{0,0.3\}$ at $M=201$; $\gamma=0$ at $M=401$ and $801$), exactly one---a
$\gamma=0$ non-learner at the densest grid, $M=801$---produces any false crossings:
four isolated single-checkpoint crossings ($0.81$--$0.88$), each rejected by the
sustained criterion; no non-learning run ever sustains. The learning runs at $M=801$
are the exception that validates
the instrument rather than breaking it: at four times the grid's budget, $4$ of $20$
$\gamma=0$ seeds genuinely solve the lookup rung in late, abrupt transitions (onsets at
steps ${\sim}32{,}000$--$38{,}400$, all in the final fifth of training; post-onset
accuracy holds above the bar on the smoothed readout, with raw checkpoint dips to
$0.66$--$0.81$), and the sustained criterion detects all four---two at zero added
latency, two attributing $7$--$8$ checkpoints later than the first single crossing
because an isolated spike precedes the ramp, the slow-ramp failure mode priced by the
confirmation-delay accounting of Section~\ref{sec:r1}. The late-emergence rate
of $0.20$ ($95\%$ CI $[0.06,0.44]$) provides the solved-reference cells that Box~1
step~7 demands, inside the TD family itself, and it sharpens the ladder's negative
result: the grid's calibration-negative reflects under-training at the studied budget,
not unlearnability.

The budget lever, however, does not generalize to a second, harder configuration. We ran a
$\gamma=0.5$ sixty-run budget grid:
$\text{steps}\in\{40{,}000,\,80{,}000\}\times T\in\{10,20,40\}\times\{$AdamW,
Muon$\}\times$ five seeds, read out with the de-noised $N=128$ evaluator so that
any emergence meets the same evidentiary bar as the rescue attempts above. The two
step budgets are four and eight times the canonical $10^{4}$ steps, counted in
optimizer steps; the checkpoint count stays at $M=201$ (evaluating every
$200$/$400$ steps), unlike the $M=801$ arm above, which scaled both axes
together. Nothing emerges: $0/60$ runs (Fig.~\ref{fig:budgetgrid}). No run reaches even the $0.7$ trainer threshold
at a single checkpoint (best transient accuracy $0.469$ across all $60\times201$
checkpoint evaluations; final accuracies $0.22$--$0.37$), so the single-crossing and
sustained criteria are silent alike.

This arm is not a clean one-factor contrast against the lookup rung, and we do not read it
as one. It differs from that rung on four axes at once: the discount ($\gamma=0.5$ against
$\gamma=0.9$), the state count ($10$ against $6$), the horizon ($T\in\{10,20,40\}$ against
$T=40$ alone), and the readout ($N=128$ evaluation sequences against $32$). Read together
with the lookup rung's $4/20$ late emergences at four times the budget, the comparison
therefore establishes a difference across a jointly varying configuration, not a
difficulty-by-budget interaction attributable to $\gamma$. What the arm does carry, and
what it is included for, is robustness of the negative itself: the best transient accuracy
anywhere in the grid is $0.469$ across all $60\times201$ checkpoint evaluations, against a
task ceiling near $0.95$, so nothing was learned in this configuration regardless of which
axis is responsible. The under-training reading of the calibration-negative is confirmed at
the lookup rung and is untested elsewhere; in this configuration the required budget, if
finite, exceeds $8\times$, and the budget lever is exhausted at this scale. A state-matched
de-confound that would isolate the discount---six states, the matched $N=32$ readout,
$T=40$, at both elevated budgets---is specified but has not been run, and we flag the
attribution as open rather than resolved.

\begin{figure}[htbp]
\centering
\figtikz[\linewidth]{E2_budget_grid}
\caption{The $\gamma=0.5$ sixty-run budget grid: per-run maximum validation
accuracy over all ${\sim}201$ checkpoints (points) and per-cell median final
accuracy (bars), by horizon $T$ and optimizer, at four and eight times the
canonical budget. No run reaches even the $0.7$ trainer threshold at a single
checkpoint (best transient $0.469$), so every detector variant is silent:
$0/60$ emergences.
(b)~Marginal maximum accuracy by horizon (all 60 runs): no horizon rescues
the task.
(c)~Crossing rate (max accuracy ${\geq}0.7$) by discount factor on the
$\gamma$ ladder: apparent emergence is a no-discount/drift phenomenon; the
$\gamma=0.5$ arm is clean at both budgets.
(d)~Best-of-grid per budget: the budget lever is exhausted ($0/60$ reach the
bar at $4\times$ or $8\times$).
Panel summaries: (a)~budget grid. (b)~horizon marginal. (c)~gamma ladder. (d)~budget exhaustion.}
\label{fig:budgetgrid}
\end{figure}

\subsection{Case study: auditing published emergence curves}\label{sec:bigbench}

The preflight computation of Box~1 applies beyond small testbeds, so we re-analysed the
public BIG-Bench score release \citep{srivastava2022bigbench}, the per-task data
underlying the emergence claims of \citet{wei2022emergent}. For the 89 multiple-choice
tasks with a defined chance baseline, evaluated on the dense BIG-G series across twelve
scales, we computed the binomial probability that a chance-level model would clear an
above-chance detector at any scale, and the expected number of such false ``emergent
tasks'' across the whole scan. A detector at chance${}+0.10$ flags 40 of the 89 tasks as
apparently emergent, yet under the all-chance null the same detector is expected to fire
on $17.7$ tasks by sampling noise alone---and the flagged count is set almost entirely
by detector leniency ($45.0$ expected false flags at chance${}+0.05$, $1.5$ at
chance${}+0.20$). A look-elsewhere correction over all $89\times12$ task-by-scale
evaluations then adjudicates the individual claims: 24 of the 40 survive
overwhelmingly---large-$N$, large-jump tasks such as \texttt{hindu\_knowledge} ($N{=}175$,
peak $0.61$ from a $0.25$ baseline, per-point exceedance $q\approx10^{-24}$) and
\texttt{date\_understanding} ($N{=}369$, $q\approx10^{-63}$)---while 9 are fragile and 7,
all small-$N$ or small-jump (e.g.\ \texttt{suicide\_risk}, $N{=}40$, jump $0.10$,
$q\approx0.10$), are unconfirmed under the scan correction: at these item counts their
jumps are indistinguishable from what a chance-level scan would produce, so they are
unpowered rather than demonstrated false. Per-task item counts and chance baselines for
the flagged claims are tabulated in Appendix~\ref{app:bigbench}.

The audit also has an error rate in the other direction, and reporting it is part of
calibrating the instrument. Eighteen of the 89 multiple-choice tasks are cited as emergent
by \citet{wei2022emergent}. The chance${}+0.10$ detector flags 40 tasks but only nine of
those eighteen, a false-negative rate of $9/18$ on precisely the claims the audit is meant
to adjudicate. Eight of the nine misses never reach chance${}+0.10$ at any scale, with peak
jumps from $0.006$ to $0.092$; the ninth, \texttt{persian\_idioms}, clears the jump at
$+0.265$ but already sits $0.205$ above chance at the smallest scale and so fails the
near-chance-early leg. Moving the bar does not repair this: at chance${}+0.05$ the detector
flags 60 tasks and still misses seven of the eighteen, and at chance${}+0.20$ it flags 18
and misses fourteen. The leniency that produces 40 flags and $17.7$ expected false ones is
thus not the same knob as the sensitivity that would recover the cited claims, and none of
the three grid settings makes the flagged set the cited set. The verdict cross-tabulation
points the same way: none of the seven unconfirmed tasks is cited, five of the 24 survivors
are, and four of the nine fragile ones are. An operational threshold-crossing rule and the
published claims are largely disjoint populations, which is a reason to report both error
directions of any such audit rather than only the false positives it was built to find.
Two scope notes.
First, the published axis is model scale (twelve points), not training checkpoints, so
the operative multiplicity is the number of tasks scanned rather than a checkpoint grid;
the correction is the same. Second, the most-cited generative emergent tasks
(e.g.\ \texttt{modified\_arithmetic}) have chance ${\approx}0$ under exact string match
and fall outside this binomial null entirely: they belong to the metric-discontinuity
mechanism of \citet{schaeffer2023mirage}, which is complementary to---not competing
with---the scan-statistic mechanism analysed here. The audit therefore corroborates
rather than debunks: it separates the emergence claims that survive a calibrated
detector from those that a preflight computation would have flagged as unpowered. The
audit script and full per-task table are included in the archived repository.

\subsection{Case study: real training trajectories (Pythia)}\label{sec:pythia}

The validations so far share one experimental pipeline, so a residual objection is that
the calibration account could be programme-local. As an external test, we ran the
detector and its correction, unchanged, on real pretraining trajectories: the public
checkpoint suites of Pythia-70M, 160M, and 410M \citep{biderman2023pythia}, scored at
twelve log-spaced public revisions (eight for 410M) spanning step $0$ to $143{,}000$, on
a twelve-task battery per model. Eight BIG-Bench multiple-choice tasks are reused from
the audit of Section~\ref{sec:bigbench}---six of its seven small-$N$
flagged-as-unpowered tasks plus two larger-$N$ tasks (\texttt{hindu\_knowledge},
$N{=}175$; \texttt{misconceptions}, $N{=}219$)---and three arithmetic prompt tasks and
LAMBADA \citep{paperno2016lambada} serve as chance-${\approx}0$ genuine-emergence probes
outside the binomial null, with a WikiText bits-per-byte control confirming smooth
training progress. The detector settings are the audit's ($\tau=\text{chance}+0.10$,
$\alpha=0.05$), with the per-curve $K$ read from Eq.~\eqref{eq:nomogram} at the measured
pre-emergence exceedance rate $\hat q$, estimated on the maximal checkpoint prefix that
ends immediately before the first sustained above-threshold block; if no such block
appears, the full curve is used, as in the drift correction of
Section~\ref{sec:ladder}.

Of the $36$ accuracy-versus-checkpoint curves, nine emerge by the end of training (last
two checkpoints above the bar) and $27$ stay at chance level (Table~\ref{tab:pythia}).
On the chance-level curves the single-crossing detector fires four times, against $5.7$
expected from the exact binomial-null computation of Box~1 step~1 over the $18$
multiple-choice chance-level curves---agreement within about one standard deviation
($1.7$); the nine arithmetic chance-level curves never leave the floor. All four fires
occur on small-$N$ tasks (\texttt{what\_is\_the\_tao}, $N{=}36$;
\texttt{common\_morpheme}, $N{=}50$, at two model sizes; \texttt{empirical\_judgments},
$N{=}99$), with measured $\hat q$ of $0.08$--$0.33$: the evaluation-size lever of
Section~\ref{sec:mnaxes} operating on public data, at exceedance rates above any
testbed cell ($\bar q=0.004$--$0.133$ across the $N$ ladder there). The
nomogram-prescribed sustained-$K$ ($K=3$--$5$ on these curves) removes all four false
fires ($0/27$; the per-curve mixture predicts $0.08$), while preserving eight of the
nine genuine emergences at zero added latency beyond the $K{-}1$ confirmation
window---including LAMBADA at all three model sizes, where a clean pre-emergence
segment ($\hat q=0$) prescribes $K{=}1$ and the correction costs nothing. The ninth is
the priced power cost of Lemma~\ref{lem:sustain}(c), not a silent failure:
\texttt{what\_is\_the\_tao} at 160M ($N{=}36$) has so noisy a pre-emergence segment
($\hat q=0.44$) that the nomogram prescribes $K{=}7$ with only three checkpoints
remaining, so its detection is censored at the end of the trace.

\begin{table}[htbp]
\centering
\caption{Pythia public-checkpoint battery: per-model detector outcomes at
$\tau=\text{chance}+0.10$ with nomogram-prescribed per-curve $K$. Parenthesised
predictions are exact binomial-null expectations over the multiple-choice chance-level
curves (the arithmetic chance-level curves have measured $\hat q=0$ and cannot fire);
the 160M censored emergence is the $K{=}7$ case discussed in the text.}
\label{tab:pythia}
\small
\begin{tabular}{@{}lccccc@{}}
\toprule
Model & Curves & Chance-level & Single fires (pred.) & Sustained-$K$ fires & Emergences caught \\
\midrule
Pythia-70M & 12 & 10 & 3 (2.7) & 0 & 2/2 \\
Pythia-160M & 12 & 9 & 1 (1.9) & 0 & 2/3 (1 censored) \\
Pythia-410M & 12 & 8 & 0 (1.1) & 0 & 4/4 \\
Total & 36 & 27 & 4 (5.7) & 0 & 8/9 \\
\bottomrule
\end{tabular}
\end{table}

Three caveats scope the replication. First, twelve checkpoints are far coarser than the
testbed's $M=201$; by Theorem~\ref{thm:single} the single-crossing false-fire
probability $1-(1-q)^{M}$ falls as the grid coarsens, so the public grid is the
lenient direction on the checkpoint axis, and this battery exercises the small-$N$
lever rather than the dense-grid lever. Second, this is one model family at three
sizes, and at $N=36$--$50$ the emerged-versus-chance split is itself weakly powered;
the unambiguous genuine emergences are LAMBADA's ($N{=}2000$).

Third, and most consequentially for how the null is read, the emerged-versus-chance
labels are instrument-derived: a curve counts as emerged when its last two checkpoints sit
above the bar, which is a member of the same threshold-crossing family whose calibration is
under test. The battery therefore validates the sustained-$K$ correction against a
partition that the uncorrected criterion supplied, and the two are not logically
independent. Two tasks make the exposure concrete by receiving different labels at
different model sizes: \texttt{what\_is\_the\_tao} is scored a false fire at 70M but a
genuine emergence at 160M and 410M, and \texttt{common\_morpheme} is a false fire at 70M
and 160M but genuine at 410M. Reclassifying the three low-size cases as genuine would move
the headline comparison from four observed fires against $5.7$ predicted to one observed,
while simultaneously removing the three highest-$\hat q$ curves from the eighteen-curve
chance-level pool and so lowering the prediction as well. We do not recompute the null on
a relabelled pool, because any defensible relabelling needs an emergence criterion
independent of the detector, which the public checkpoint grid does not supply; the point of
reporting the size-inconsistent cases is that the $4$-versus-$5.7$ agreement is contingent
on a detector-derived split rather than on ground truth. Within that scope, both
halves of the account---the predicted single-crossing false-fire rate and a
sustained-$K$ correction that suppresses noise fires while sparing genuine
transitions---operate as predicted on pretraining checkpoints produced entirely
outside this programme's pipeline.

\subsection{Case study: data-presentation order does not compress the emergence delay}\label{sec:r1}

As a case study for the floor-censoring and learnable-target disciplines
(Section~\ref{sec:disciplines}), we ask whether data-presentation order can compress a
genuine emergence delay.
On the in-context copy/induction family---the one task in our suite that is learnable,
non-Fourier, and shows a measurable emergence delay---we sweep four ordering policies
(i.i.d., easy$\to$hard, hard-first, structured) $\times$ $\{$AdamW, SGD with momentum
(SGDM)$\}$ $\times$ 5 seeds.
The induction circuit emerges at 1750--1950 steps regardless of ordering; the i.i.d.\
(no-curriculum) baseline is the fastest (mean 1780) and the easy$\to$hard
curriculum is mildly but significantly slower (mean 1900; pooled Welch
$t = -2.83$, $p = 0.011$, $n = 10$; the AdamW-only contrast is borderline at $t = -2.23$,
$p = 0.058$, $n = 5$). These ordering-policy contrasts are exploratory descriptive checks,
reported unadjusted rather than as a separate corrected confirmatory family. The honest
statement for this copy/induction family is therefore ``no
ordering helps and easy$\to$hard mildly hurts,'' not a flat zero-effect null: the
emergence delay is set by the task, and ordering does not substitute for target structure
(Fig.~\ref{fig:curriculum}; Table~\ref{tab:curriculum}). (The Walsh and modular-addition families did not yield an
adjudicable readout---a testbed boundary we document rather than over-read; the Walsh
target is unlearnable in this pipeline (Fig.~\ref{fig:walsh-control}), which is itself one
of the calibration lessons.)

The induction circuit itself is a clean positive emergence---a useful contrast to the
in-context TD non-emergence in Section~\ref{sec:r2}: across context lengths the circuit appears within
a few hundred steps and every cell reaches a high in-context-learning score
(Fig.~\ref{fig:induction-emergence}; these context-length steps use a different readout and
task setup from the ${\approx}1850$-step ordering measurement above and are not directly
comparable), so the delay we measure is a property of a circuit
that genuinely forms, not of a detector firing on noise. The positive-control grid here
adds a third optimizer, Muon, to the
$\{$AdamW, SGDM$\}$ pair used in the ordering sweep above, giving the
$\{$AdamW, SGDM, Muon$\}$ set referenced in the figures and tables.

\paragraph{The sustained criterion validated on real transitions.} The noise suppression
of Section~\ref{sec:r2} would be worthless if sustaining also suppressed or delayed
genuine emergence. It does not. Applying both detectors at the same $\tau=0.8$ bar to all
90 induction runs (3 optimizers $\times$ 3 context lengths $\times$ 10 seeds; evaluation
every 100 steps), every run truly emerges (final repeat-token accuracy $\geq0.9$) and both
detectors fire in $90/90$ runs. Critically, sustaining is free on this side: the
sustained-2 detection step equals the single-crossing step in $90/90$ runs---zero
latency at the evaluation cadence---because a genuine transition is sticky (once the
signal crosses the bar it stays above it), whereas a noise crossing is isolated by
construction. Even sustained-3 has median zero latency; eleven of the $90$ runs pay a
nonzero attribution cost of three to five evaluation intervals, two of them the
maximum five. Requiring $K=2$ consecutive exceedances therefore buys the noise-side
false-positive suppression quantified in Corollary~\ref{cor:numbers} at zero measured
detection cost on a genuinely emerging family.

Two clarifications keep this claim honest. First, ``zero latency'' is an
onset-attribution statement: the sustained detector reports the first checkpoint of the
qualifying run, and by construction it cannot confirm an emergence until $K-1$
further checkpoints have been observed, so the confirmation delay is always
$(K{-}1)\times$ the evaluation cadence. Second, the transitions above are sharp, which
is the criterion's most favorable case, so we engineered toward the unfavorable one:
rerunning the copy/induction family at $3.3\times$ and $6.7\times$ reduced learning
rates with a $4\times$ finer evaluation cadence (25 steps). All ten seeds still emerge,
${\sim}3.7\times$ later than at the default rate (onsets 6525--6850), and the
transitions remain overwhelmingly sharp: accuracy climbs from $0.3$ to above $0.8$
within one to three 25-step checkpoints in nine of ten runs, while the tenth produces
the sought gradual ramp---eleven checkpoints (${\sim}275$ steps). In all ten runs,
including on that ramp, single-crossing, sustained-2, and sustained-3 attribute the
identical onset, because the ramps are monotone once under way---no pre-ramp noise
spike preceded any of them. Across the $104$ genuine transitions observed in this
paper's testbeds ($90$ induction, $10$ reduced-lr copy, $4$ late TD emergences),
sustaining adds attribution latency in exactly two---the two late TD emergences whose
ramps are preceded by an early isolated spike ($7$--$8$ checkpoints each;
Section~\ref{sec:mnaxes})---and is latency-free in the other $102$; the failure mode
(an early isolated spike on a slow ramp) is thus observed, rare, and exactly what
the confirmation-delay accounting above prices.

\begin{figure}[t]
\centering
\figtikz[\linewidth]{E2_curriculum}
\caption{Emergence step by data-ordering policy (copy/induction family; mean $\pm$ s.d.\
over 5 seeds per optimizer). (a)~Across the four ordering policies the circuit emerges at
1750--1950 steps; the i.i.d.\ baseline is the fastest (mean 1780 steps) and easy$\to$hard
is the slowest (mean 1900).
(b)~Per-seed emergence steps by ordering: the null is seed-consistent.
(c)~Delta versus each optimizer's own iid baseline: no ordering compresses
the delay.
(d)~Family robustness on copy versus mod-add (120-run grid): the
no-compression verdict holds on both families; Walsh is not in this panel.
Panel summaries: (a)~emergence step by ordering and optimizer. (b)~per-seed emergence steps. (c)~delta vs iid baseline. (d)~copy vs mod-add robustness.}
\label{fig:curriculum}
\end{figure}

\begin{figure}[htbp]
\centering
\figtikz[\linewidth]{E2_induction}
\caption{The induction-circuit emergence is a clean positive control (contrast the
non-emergence of Section~\ref{sec:r2}). (a)~Median emergence step (interquartile range over 10 seeds) versus context
length for three optimizers; the circuit forms within a few hundred steps in every cell, with
plain-momentum SGD the slowest. The emergence step is read on a 100-step evaluation grid, so
the fastest cells are grid-limited at the floor and the AdamW/Muon difference is not resolved
here. These context-length emergence steps (a few hundred) use a different readout and task
setup from the ${\approx}1850$-step copy/induction emergence of Section~\ref{sec:r1}
(Fig.~\ref{fig:curriculum}, threshold $0.5$); the two are not directly comparable.
(b)~Final in-context-learning score per optimizer (each point a seed): all cells sit
well above the $0.5$ reference, confirming the circuit genuinely forms rather than being a
detector artifact.
(c)~Sustained-$K$ latency audit: the sustained-$K{=}2$ crossing adds at most
one checkpoint over the single crossing per run (median zero), so the
correction costs no latency on genuine emergences.
(d)~Per-seed onset steps by context length: the raw distributions behind the
median ordering.
Panel summaries: (a)~emergence step vs context length. (b)~final ICL score by optimizer. (c)~sustained-$K$ latency. (d)~per-seed onset steps.}
\label{fig:induction-emergence}
\end{figure}

\begin{table}[htbp]
\centering
\caption{Emergence step per ordering policy and optimizer (copy/induction family, $n=5$ seeds each).
Median emergence step at threshold $0.5$; cell means span 1780--1900 and cell medians 1750--1950 across orderings.
Pooled Welch $t$-test (i.i.d.\ vs.\ easy$\to$hard, $n=10$ each): $t=-2.83$, $p=0.011$.}
\label{tab:curriculum}
\small
\begin{tabular}{llcc}
\toprule
Ordering & Optimizer & Median step & Mean step \\
\midrule
i.i.d.\       & AdamW & 1750 & 1780 \\
i.i.d.\       & SGDM  & 1850 & 1830 \\
easy$\to$hard & AdamW & 1950 & 1900 \\
easy$\to$hard & SGDM  & 1850 & 1880 \\
hard-first    & AdamW & 1850 & 1840 \\
hard-first    & SGDM  & 1850 & 1850 \\
structured    & AdamW & 1850 & 1830 \\
structured    & SGDM  & 1850 & 1850 \\
\bottomrule
\end{tabular}
\end{table}

\subsection{Case study: early trajectories are a configuration fingerprint, not a run predictor}\label{sec:r3}

As a case study for the nuisance-axis-freezing discipline, we ask whether early-trajectory
predictiveness survives once the optimizer axis is frozen.
Re-analysing 364 existing runs across 11 closed experiment namespaces (252 grok events each
labelled with a route---a discrete outcome-class assigned to the run and used here as the
prediction target; zero new training), we test three early-prediction laws (a seed-level
delay-signal increment, a route-prediction area under the ROC curve (AUROC), and a config-baseline error
comparison)---whose analysis plan (the three laws plus baselines) was pre-registered before
the data were inspected, although the runs themselves are pre-existing rather than collected
prospectively---on this first-window trajectory against a configuration-identity baseline,
with group-aware cross-validation and label-permutation nulls. All three land on their pre-registered
clean-negative branches: the seed-level signal increment for delay is statistically
detectable but practically nil (Spearman $\rho \approx 0.077 / 0.014$) and the signal
model's error is worse than a predict-the-config baseline. The quantified warning
for the timing-prediction literature: the pooled route AUROC of $0.815$ collapses toward
chance ($0.47$--$0.57$, not distinguishable from chance at this sample size)
within AdamW---optimizer identity masquerades as trajectory
predictiveness unless the optimizer axis is frozen. Reported predictability that pools
across optimizers can be configuration recognition, not run prediction.
The full window~$\times$~tier breakdown is in Fig.~\ref{fig:route-matrix}, and the
collapsed summary is in Fig.~\ref{fig:specroute}. This recasts early-prediction claims
\citep{notsawo2023predicting, xu2026earlywarning} as requiring a frozen nuisance axis
before the predictiveness can be attributed to the run.

\begin{figure}[htbp]
\centering
\figtikz[0.7\linewidth]{E2_routemat}
\caption{Tier-B route-prediction AUROC at both early windows, pooled across optimizers
(bars); the dotted line marks chance. Tier~B has no estimable within-AdamW cell
(insufficient within-optimizer label balance), as noted in the panel.
(b)~Tier-A reference: the pooled-versus-within-AdamW collapse (pooled $0.815$/$0.67$
dropping to $0.57$/$0.47$ within a single optimizer, across the 200-step and mem-onset
windows; the four bars of the main-text spec-route figure, restated).
(c)~The optimizer-identity component per window (pooled minus within-AdamW,
tier A).
(d)~Variance decomposition at the mem-onset window: config share dominates in
every optimizer---the early window is a configuration fingerprint, not a run
fingerprint.
Panel summaries: (a)~tier-B AUROC. (b)~within-AdamW collapse. (c)~optimizer component. (d)~variance decomposition.}
\label{fig:route-matrix}
\end{figure}


\begin{figure}[htbp]
\centering
\figtikz[\linewidth]{E2_specroute}
\caption{Four diagnostics for the early-prediction collapse, panels (a)--(d).
(a)~Signal increment is not stable: the change in grok-prediction AUROC from adding the
trajectory signal on top of the config-only baseline, for the config, seed, and task
predictors (the mem-onset window W-mem and the 200-step window W-200). Increments are at most $+0.07$ and the task
predictor moves $-0.38$, so the signal adds no reliable predictiveness over config.
(b)~Pooled route AUROC is optimizer identity: pooled route-prediction AUROC is high
($0.815$ at the 200-step window, $0.674$ at mem-onset) but collapses to chance within
AdamW ($0.57$/$0.47$, below the $0.5$ dotted line), exposing that pooled predictiveness
is configuration recognition, not run-level forecasting.
(c)~Config baseline wins: leave-one-out mean absolute error (MAE) of log-delay residual for the signal model
versus the predict-the-config baseline; the signal model's error ($0.51$/$0.52$) exceeds
the config baseline ($0.33$/$0.43$) at both windows.
(d)~Config dominates optimizer structure: out-of-fold variance-share decomposition of
log-delay per optimizer (config, signal, residual). AdamW and Muon are config-dominated;
SGDM carries no recoverable structure (see below).
In the variance-share decomposition, shares are clamped at zero and need not
sum to one: for SGDM both the config and signal out-of-fold $R^2$ are negative
(predictors are worse than the group mean), so config and signal collapse to zero and the
residual share exceeds one ($1.26$), indicating that no fitted component explains its
delay variance.
Panel summaries: (a)~$\Delta$ AUROC. (b)~pooled vs within-AdamW AUROC. (c)~LOO MAE. (d)~variance share.}
\label{fig:specroute}
\end{figure}

\section{Reconciliation and scope}\label{sec:discussion}

\paragraph{The shared root: measurement before interpretation.}
All three case studies trace to the same preventable failure mode: interpreting a measurement
output before verifying that the measurement protocol can distinguish a real signal from
its null.
A single-crossing detector applied to a stationary accuracy series
will reliably fire (Section~\ref{sec:theorem}); an unlearnable target will reliably fail any ordering
intervention; a predictor trained without freezing the optimizer axis will reliably
absorb the optimizer label.
These are not unusual mistakes---they are precisely the traps that arise when
experiments are cheap enough to run at scale but too fast to pre-calibrate carefully.

\paragraph{Relation to the large-model emergence debate.}
Our findings mirror, at small scale, the concern raised by \citet{schaeffer2023mirage}
for large models: apparent phase transitions can be artefacts of the measurement
protocol rather than genuine discontinuities in the underlying computation.
The mechanism differs---metric nonlinearity for large models, noisy single-crossing
detectors and unlearnable targets for small models---but the corrective logic is the same.
The single-crossing detector is the small-scale counterpart of a discontinuous metric: it
maps a noisy accuracy series onto a binary ``emerged / not'' label by asking whether the
series ever crosses a threshold.
Applied to our in-context TD grid (45 runs, median final accuracy $0.562$, against a
constant-predictor floor of $0.390$ and a detection bar of $0.8$)
at the illustrative $\tau=0.8$ bar, it manufactures apparent emergence in 26/45 runs, whereas the sustained-2 criterion---the
counterpart of a continuous or de-noised readout---removes every apparent emergence (0/45
survive).
Theorem~\ref{thm:single} and Lemma~\ref{lem:sustain} give the closed-form, parameter-free account: with $N=32$ evaluation
sequences, $M=201$ evaluation points, and stationary baseline $p=0.562$, a single crossing
fires with probability $0.421$ per run without any transition while a sustained-2 crossing fires with
probability $0.00147$---a suppression of more than two orders of magnitude (under the
stated association hypothesis). At the pre-registered $\tau=0.7$ bar the corresponding chain is
$42/45\to1/45$ under the mixture-prescribed $K=5$ (Section~\ref{sec:theorem}).
The instrument (detector choice) manufactures the phenomenon (emergence rate), in direct
analogy to \citeauthor{schaeffer2023mirage}'s finding that the metric manufactures the
phenomenon in large-model scaling. This is a reconciliation rather than a new phenomenon: in
the genre of \citet{oneto2026}, the apparent effect is accounted for by a known
mechanism---here, scan-statistic false positives.

\citet{lee2023icrl} provide the positive existence proof that in-context RL can
emerge: their supervised-pretraining recipe on diverse RL trajectories yields near-optimal
in-context RL agents at test time.
Our TD testbed does not refute this; it is a calibration failure on a harder, single-task
online regime.
The honest framing is: \citeauthor{lee2023icrl} establish that in-context TD emergence is
achievable in principle, which is precisely why our recalibration recipe (solved-reference
cell, de-noised readout, sustained criterion) matters: it sets the minimum calibration bar
an online single-task testbed must clear before emergence can be claimed, and without such a
gate there is no way to tell a broken testbed from a genuine negative.

\paragraph{Curriculum learning and emergence timing.}
The curriculum learning literature \citep{bengio2009curriculum, soviany2022curriculum, wu2021curricula}
generally shows benefits for asymptotic performance or sample efficiency, not for
shifting the onset of phase-transition-like learning.
Our ordering case study (Section~\ref{sec:r1}) is consistent with this: if the emergence delay is set by the task
structure (the induction circuit must be assembled regardless of input order), then
presenting examples in a ``better'' order cannot materially accelerate that assembly.
Easy$\to$hard ordering mildly lengthens the delay, plausibly because the early
easy-example phase trains representations that must be partially unlearned when harder
examples arrive.
Curricula can help or hurt depending on
whether the hard-example representations subsume the easy-example ones.

\paragraph{Pre-registration as a methods virtue.}
Following the methodological prescription of \citet{pineau2021reproducibility},
all three case studies had their analysis plans---thresholds, comparison arms, and
outcome branches---fixed before analysis (for the re-analysis of
Section~\ref{sec:r3}, before the pre-existing runs were inspected).
Pre-registration is especially useful for calibration work because the
``interesting'' finding---a null that rules out a hypothesis---requires the same
evidential discipline as a positive finding: without it, a researcher who runs four ordering
policies and finds a borderline-significant slowdown could plausibly report that as a
positive curriculum effect. We treat it as one methodological safeguard among several, not
as the identity of the paper.

\paragraph{Scope and forward directions.}
The validation here is specific to a single programme of tiny-transformer experiments; the
theorem is general but the empirical anchoring is local.
The calibration disciplines below (Section~\ref{sec:disciplines}) follow from the reduction
as transferable recommendations, not as universal theorems.
Extending the in-context TD testbed to a regime where emergence is achievable would
convert the current calibration vignette into an adjudicable science question; the
budget and discount levers are now answered in the negative at this scale
(Section~\ref{sec:mnaxes}), leaving pretraining diversity---for instance, via the
approach of \citet{lee2023icrl}---as the untested lever.
Similarly, the early-prediction collapse (Section~\ref{sec:r3}) motivates a cleaner experimental
design: fix the optimizer and vary only seed and data, then ask whether within-seed
early trajectories predict final performance.
The sustained-crossing criterion of Section~\ref{sec:theorem} is immediately applicable to
any binary-accuracy emergence study as a drop-in replacement for single-crossing
detectors, with an observed ${\sim}42\times$ reduction in false positives at the
pre-registered bar of the testbed studied here.

\subsection{The calibration disciplines}\label{sec:disciplines}

The three case studies share a small set of preventable failure modes, which the reduction
of Section~\ref{sec:theorem} implies as calibration disciplines for small-scale emergence
studies:

\begin{enumerate}
\item \textbf{Sustained vs single-crossing criteria.} A single threshold crossing is
  noise; require the signal to persist for $K$ consecutive evaluations. Theorem~\ref{thm:single}
  and Lemma~\ref{lem:sustain} quantify the $q^{-(K-1)}$ suppression, and
  Eq.~\eqref{eq:nomogram} sets $K$ for a target budget (Section~\ref{sec:r2}, and cf.\
  mode-connectivity and norm-clock studies).
\item \textbf{Floor-censoring control.} Coarse evaluation grids censor fast transitions at
  a floor; a fine-evaluation arm is needed before reading ``instant'' or ``zero-variance''
  (Section~\ref{sec:r1}, mod-add arm; cf.\ Edge-of-Stability \citep{cohen2021eos} and induction-emergence studies).
\item \textbf{Learnable-vs-unlearnable target check.} An intervention that fails on an
  unlearnable target (e.g.\ the Walsh pure-degree target) is uninformative; verify the
  target is learnable in the pipeline first (Section~\ref{sec:r1}).
\item \textbf{Freeze nuisance axes before claiming predictiveness.} Predictiveness pooled
  over a nuisance axis (here the optimizer) can be recognition of that axis; freeze it
  (Section~\ref{sec:r3}).
\item \textbf{A recurring negative anchor.} Plain-momentum SGD repeatedly serves as a clean
  negative anchor (it fails to train several of these tasks), useful for calibrating ``did
  anything happen at all.''
\end{enumerate}
Table~\ref{tab:evidence-ledger} summarises each case study with its diagnostic, evidence scope, and what further work would be needed before a stronger claim.

\begin{table}[htbp]
\centering
\caption{Validation summary across the case studies. The final column
marks what would be needed before making a stronger positive/negative claim.}
\label{tab:evidence-ledger}
\scriptsize
\begin{tabular}{@{}>{\raggedright\arraybackslash}p{0.23\linewidth}>{\raggedright\arraybackslash}p{0.24\linewidth}>{\raggedright\arraybackslash}p{0.18\linewidth}>{\raggedright\arraybackslash}p{0.25\linewidth}@{}}
\toprule
Claim & Diagnostic & Evidence & Boundary \\
\midrule
Order does not compress delay & Fixed-target order arms and solved-reference check & Copy/induction ordering grid & Copy/induction family only \\
In-context TD remains a calibration failure in the trained grid & Sustained threshold plus de-noised detector & TD audit and false-positive case study & Calibration negative; not a theorem that TD cannot emerge \\
Early trajectories confound optimizer identity & Pooled vs within-optimizer AUROC & Trajectory re-analysis & Warning is local to this corpus/optimizer mix \\
Detector must reject noise & Single-crossing vs sustained emergence & False-positive trajectory audit plus 171 failed trainable-control cells (rescue, ladder, $\gamma=0.5$ budget grid) & Oracle validates labels; trained positive remains absent at up to $8\times$ budget \\
Calibration transfers off-programme & Audit-convention detector plus nomogram-prescribed sustained-$K$ on public checkpoints & Pythia $70$M--$410$M, $36$ real pretraining curves ($4\to0$ false fires, $8/9$ emergences kept) & Coarse grid ($M{=}12$/$8$), one model family; small-$N$ tasks drive the fires and the one censored detection \\
Negative anchor is useful & SGDM failure across related testbeds & Cross-paper optimizer floor evidence & Anchor diagnoses training failure, not mechanism by itself \\
\bottomrule
\end{tabular}
\end{table}

\subsection{Limitations}\label{sec:limitations}

These are toy-scale results on one programme's testbeds; the negatives follow analysis
plans fixed in advance and are clean but individually modest. The Pythia battery
(Section~\ref{sec:pythia}) shows the detector account itself transfers off-programme,
but that replication is scoped in its own right: twelve (eight for 410M) log-spaced
public checkpoints per curve rather than $201$---by Theorem~\ref{thm:single} the
direction that reduces single-crossing false fires, so it tests the small-$N$ lever
rather than the dense-grid lever---one model family at three sizes, and
multiple-choice item counts small enough ($N{=}36$--$50$ on several tasks) that the
emerged-versus-chance split is itself weakly powered. The ordering adjudication is on the copy/induction family
only (other families were testbed-bounded). The in-context TD grid is a calibration-failure vignette---it does
not establish that in-context TD cannot emerge, only that the trained grid cannot
adjudicate it despite a solved-reference oracle validating the evaluator; indeed, at
four times the budget the lookup rung does produce genuine late emergences ($4/20$
seeds; Section~\ref{sec:mnaxes}). The same lever fails in a second, harder configuration:
the $\gamma=0.5$ sixty-run budget grid stays negative at eight times the budget ($0/60$;
Section~\ref{sec:mnaxes}). That arm varies the state count, horizon set, and readout size
alongside the discount, so it bounds the negative rather than isolating which axis
resists the budget; the two results are not in contradiction.
Additional isolated trainable-control attempts made the TD task easier
($n_{\mathrm{states}}\in\{2,4\}$, $T\in\{4,6\}$; 26 trained cells, across two attack rounds, spanning AdamW/Muon,
ultra-easy AdamW, and a wider 897k-parameter AdamW rescue). None obtained a sustained
$0.8$ run under the detector (Fig.~\ref{fig:positive-rescue}). The wider rescue was informative but still not a positive
control: all three cells sustained $0.7$ yet peaked only at $0.754$ and ended near


\begin{figure}[htbp]
\centering
\figtikz[\linewidth]{E2_positive_rescue}
\caption{Trainable-positive rescue: near-positive but not a positive control.
(a)~Best max validation accuracy per rescue family, pooling both attack rounds
(main: $n=14$ cells; ultra-easy: $n=9$; wide 897k-parameter AdamW: $n=3$); the wide
rescue improves toward $0.754$ but does not
cross the $0.8$ positive-control threshold.
(b)~95\% one-sided upper bound on the success probability given $0/26$ successes
($\hat{p}_{\mathrm{upper}}=0.109$): this is an upper bound under this sampled rescue set,
not an impossibility proof. It also assumes independent cells; the 26 cells share one
pipeline across three arms, so the cluster-aware reading is weaker
(Section~\ref{sec:limitations}).
This is why Section~\ref{sec:r2} is framed as a non-adjudicable TD testbed rather than evidence that
TD emergence cannot occur.
(c)~Near-positive ($0.7$) counts per family: the wide family comes closest,
but near-positive is not a positive control.
(d)~Pool provenance: the 26 rescue cells pooled across both attack rounds
with zero successes (95\% upper bound $0.109$ on the per-cell rate).
Panel summaries: (a)~best rescue accuracy. (b)~zero-success bound. (c)~near-positive cells. (d)~pool provenance.}
\label{fig:positive-rescue}
\end{figure}
$0.69$--$0.71$ (Fig.~\ref{fig:positive-rescue}). The discount ladder of
Section~\ref{sec:ladder} extends this search to the family's easiest rung ($V(s)=r(s)$,
an in-context lookup) and still finds no trained positive at the standard budget, and
the $\gamma=0.5$ sixty-run budget grid (Section~\ref{sec:mnaxes}) extends it along the
budget axis. Across all 171 trainable-control cells---the 26 rescue cells, the 85
standard-budget ladder runs at $\gamma>0$, and the 60 budget-grid runs---zero genuine
sustained emergences were observed. These cells are not independent draws: all 171 share
one pipeline, architecture, and tokenizer, and they group into nine experimental arms
(three rescue families, four ladder rungs, two budget tiers). Under the optimistic
independence assumption the $0/171$ count gives a 95\% upper bound of about $0.02$ on the
per-cell success probability of this family of rescues; under the pessimistic extreme in
which each arm contributes one effective trial, the $0/9$ count gives about $0.28$. A
hierarchical (Beta-Binomial) treatment interpolates between these endpoints, but with
zero successes the within-arm dependence is unidentified from the data, so we report the
result descriptively---$0$ successes in $171$ cells across $9$ arms---with the honest
per-cell bound lying between $0.02$ and $0.28$ depending on how strongly the shared
pipeline correlates outcomes. (The $\gamma=0$ lookup rung is excluded from
the pool by competence class, not outcome: it reduces to in-context retrieval rather
than value bootstrapping, and it is where the family's only trained positives
live---the $4/20$ at quadrupled budget quoted above.) Thus the exact Bellman cell remains an oracle/solved-reference control that
validates the evaluator; the trained positive controls for the detector are supplied by
the induction family (Section~\ref{sec:r1}), where all 90 genuine emergences are
detected at zero sustained-criterion latency, and by the lookup rung's late emergence
at the quadrupled budget (Section~\ref{sec:mnaxes}). We therefore further narrow the
claim: the budget lever is exhausted at this scale (silent at eight times the
canonical budget on the $\gamma=0.5$ configuration), and this setup is not a reliable
trainable TD-emergence testbed under any budget probed; whether other levers,
pretraining diversity above all, unlock it remains open.
The early-prediction case study is a post-hoc re-analysis at a fixed corpus scale; the AUROC-collapse warning is
specific to corpora where the optimizer axis is varied.
The within-AdamW collapse (W-200 AUROC $0.57$, $p\approx0.12$, $n=72$) is power-bounded:
a true AUROC around $0.6$ cannot be excluded at this sample size, so the collapse is a
quantified warning, not a proof that within-AdamW predictiveness is exactly at chance.
Finally, the false-positive account itself uses classical run/scan statistics rather than a
new inequality; its novelty is the reduction and the calibrated instrument. The
illustrative-bar suppression figure of Corollary~\ref{cor:numbers} assumes near-independence, which is now measured rather than pending
(median detrended lag-1 autocorrelation $+0.03$, $41/45$ runs inside the white-noise
band, with lags $2$--$5$ equally quiet; a deliberately pessimistic dependence model still
leaves ${\approx}24\times$, Section~\ref{sec:theorem}).

\section{Conclusion}\label{sec:conclusion}

We have reduced the emergence-timing readout used in small-scale testbeds to a classical
run/scan-statistics false-positive problem, and calibrated a sustained-crossing correction
with an explicit nomogram.
The central message is not that emergence is rare or unreal, but that small-scale
testbeds require the same measurement discipline as large-scale ones: a solved-reference
cell to validate the evaluator, a sustained threshold to suppress noise-driven detections,
a learnable-target check before any intervention is tested, and a frozen nuisance axis
before any predictiveness claim is made.
The false-positive theorem (Section~\ref{sec:theorem}) shows that the suppression
delivered by the sustained-crossing criterion---an observed ${\sim}42\times$ at the
pre-registered bar---is a direct consequence of the evaluation
parameters (under the stated association hypothesis).
The mathematics is classical; what is new is the reduction of small-scale emergence timing
to it, the sustained-$K$ power/latency analysis, and the calibrated instrument.
We hope these diagnostics serve as reusable infrastructure for the growing community
studying phase-transition-like learning at small scale.

\appendix
\setcounter{figure}{0}\renewcommand{\thefigure}{A\arabic{figure}}
\setcounter{table}{0}\renewcommand{\thetable}{A\arabic{table}}
% Keep hyperref anchors unique after the counter reset (avoids duplicate-destination warnings)
\renewcommand{\theHfigure}{A\arabic{figure}}
\renewcommand{\theHtable}{A\arabic{table}}
\section{Supplementary run-grids and robustness checks}\label{app:grid}

These supplementary figures back the three main-text case studies, in the order they are
first cited: a fit-sanity check for the in-context TD testbed
(Fig.~\ref{fig:supervised-fit}), the trainable-positive rescue audit
(Fig.~\ref{fig:positive-rescue}), and the consolidated run-grids behind all
three nulls (Fig.~\ref{fig:grid}).
Two further robustness figures illustrate calibration disciplines directly: the
unlearnable-target check of Section~\ref{sec:disciplines} (Fig.~\ref{fig:walsh-control}) and
the tier-B route-prediction AUROC at both early windows
(Fig.~\ref{fig:route-matrix}), complementing the tier-A collapse of Fig.~\ref{fig:specroute}.



\begin{figure}[htbp]
\centering
\figtikz[\linewidth]{E2_walsh}
\caption{Learnable-vs-unlearnable target check (calibration discipline~3). (a)~On a single
non-Fourier target, final accuracy stays pinned at the floor ($\approx0.04$--$0.05$) across
a $20\times$ range of total training budget ($2{,}400$ to $48{,}000$ steps; markers are
per-seed runs, diamonds the per-budget median, all AdamW), far below the learnable copy
reference ($1.00$, dashed line). An intervention tested on such a target is uninformative,
which is why the ordering sweep of Section~\ref{sec:r1} is adjudicated on the copy/induction family
rather than this one.
(b)~The floor is ordering-robust: the non-Fourier target stays at the floor
under every data-ordering policy.
(c)~Budget-response per target (medians): the learnable target improves with
budget; the non-Fourier target does not (${\sim}20\times$ span).
(d)~The learnable-versus-floor accuracy margin by budget: the discriminative
margin of the discipline-3 control.}
\label{fig:walsh-control}
\end{figure}

The consolidated run-grids (Fig.~\ref{fig:grid}) rest on $\approx$255 runs across the
three nulls, in one panel-per-result view. Panel~(a): across the ordering sweep (120
runs; three task families $\times$ four orderings $\times$ two optimizers $\times$
five seeds) the copy/induction circuit emerges flat at $\sim$1850 steps regardless
of ordering, with i.i.d.\ the fastest---mod-add floors at the first evaluation
(no headroom to compress) and the Walsh family does not emerge within budget.
Panel~(b): the in-context-TD grid (45 runs) sits at $\approx$chance for every
optimizer $\times$ horizon cell; Table~\ref{tab:tdcalib} shows that de-noised, sustained
threshold criteria remove the single-spike positives.
Panel~(c): the route-prediction AUROC is high when pooled across optimizers
($0.815$ at the 200-step window) but collapses to chance within AdamW
(0.57/0.47), exposing the early window as a configuration fingerprint.





\begin{figure}[htbp]
\centering
\includegraphics[width=\linewidth]{E2_grid.pdf}
\caption{Run-grids behind the three pre-specified negatives ($\approx$255 runs total).
(a)~Data-ordering sweep (120 runs; copy/induction family flat at $\sim$1850 steps with
i.i.d.\ the fastest; mod-add floors at the first evaluation; walsh does not emerge within
the training budget --- note log-$y$ axis).
(b)~In-context TD grid (45 runs; 3 optimizers $\times$ $T\in\{10,20,40\}$ $\times$ 5
seeds); each cell shows median accuracy / runs whose accuracy ever crosses $0.8$ at a
single checkpoint, of $n=5$---$26/45$ runs show such a transient single crossing,
none is sustained ($0/45$), and the grid is $\approx$chance everywhere.
(c)~Early-window route-prediction AUROC (90+ route-labelled runs, subset of 364 total);
pooled signal ($0.815$) collapses to
chance within AdamW ($0.57$/$0.47$), confirming that the early window is a configuration
fingerprint.}
\label{fig:grid}
\end{figure}

\section{BIG-G audit: per-task item counts and chance baselines}\label{app:bigbench}
\setcounter{table}{0}\renewcommand{\thetable}{B\arabic{table}}
\renewcommand{\theHtable}{B\arabic{table}}

The look-elsewhere adjudication of Section~\ref{sec:bigbench} is driven by each task's
item count $N$ and chance baseline, so we document both. Across the 89 audited
multiple-choice tasks, $N$ ranges from $8$ to $70{,}000$ (median $175$) and the chance
baseline from $0.012$ to $0.5$. The verdict classes separate cleanly by power: the 24
surviving claims have median $N=188$ (range $8$--$54{,}668$; the $N{=}8$ survivor clears
on a $0.75$ jump), the 9 fragile claims median $N=100$ ($32$--$219$), and the 7
unconfirmed claims median $N=50$ ($36$--$99$)---the flagged-but-unconfirmed tasks sit at
the small-$N$, small-jump corner where a chance-level scan is expected to produce
equally extreme flags. The seven unconfirmed tasks are listed in the appendix table below;
$E_{\mathrm{task}}$ is the expected number of chance-level tasks in the 89-task scan
showing a jump at least as extreme, so values $\gg1$ mean the observed jump is
unremarkable under the all-chance null. None of the seven is among the headline tasks
cited by \citet{wei2022emergent}; the audit's verdict for them is ``unpowered under
scan correction,'' not ``demonstrated false.'' The full 89-row table is in the archived
repository.

The detector's false negatives on the cited claims (Section~\ref{sec:bigbench}) come from
the same inputs. Of the eighteen cited-emergent multiple-choice tasks, the nine that the
chance${}+0.10$ detector does not flag are listed in
Table~\ref{tab:bigbench-falsenegative}. Eight fail the detector's jump leg, never reaching
chance${}+0.10$ at any of the twelve scales; \texttt{persian\_idioms} fails the
near-chance-early leg instead, since it already scores $0.455$ at the smallest scale
against a chance baseline of $0.250$. Item counts run from 32 to 323 (median 60), so most
sit in the same small-$N$ corner as the seven unconfirmed flags of
Table~\ref{tab:bigbench-unpowered}: a scan correction has little power to confirm these
claims, and the detector has little power to find them.

\begin{table}[htbp]
\centering
\caption{The nine cited-emergent BIG-G multiple-choice tasks that the audit detector does
not flag. Jump = peak accuracy minus chance; early = accuracy at the smallest scale.
A task is flagged only if jump ${\geq}0.10$ and early ${\leq}$ chance${}+0.05$.}
\label{tab:bigbench-falsenegative}
\footnotesize
\begin{tabular}{>{\raggedright\arraybackslash}p{0.36\linewidth}ccccl}
\toprule
Task & $N$ & Chance & Peak & Jump & Leg failed \\
\midrule
\texttt{persian\_idioms}         & 66  & 0.250 & 0.515 & 0.265 & early \\
\texttt{english\_proverbs}       & 34  & 0.290 & 0.382 & 0.092 & jump \\
\texttt{international\_phonetic\_}\newline\texttt{alphabet\_nli}
                                 & 126 & 0.333 & 0.397 & 0.064 & jump \\
\texttt{analogical\_similarity}  & 323 & 0.143 & 0.189 & 0.046 & jump \\
\texttt{crass\_ai}               & 44  & 0.250 & 0.295 & 0.045 & jump \\
\texttt{odd\_one\_out}           & 86  & 0.206 & 0.244 & 0.038 & jump \\
\texttt{code\_line\_description} & 60  & 0.248 & 0.283 & 0.035 & jump \\
\texttt{known\_unknowns}         & 46  & 0.500 & 0.522 & 0.022 & jump \\
\texttt{gre\_reading\_comprehension} & 32 & 0.284 & 0.290 & 0.006 & jump \\
\bottomrule
\end{tabular}
\end{table}

\begin{table}[htbp]
\centering
\caption{The seven flagged-but-unconfirmed BIG-G emergence claims (detector: accuracy
$\geq$ chance${}+0.10$ at some scale, near-chance early). $N$ = items; jump = peak
accuracy minus chance; $E_{\mathrm{task}}$ = expected number of equally extreme flags
across the 89-task scan under the all-chance null.}
\label{tab:bigbench-unpowered}
\small
\begin{tabular}{lccccc}
\toprule
Task & $N$ & Chance & Peak & Jump & $E_{\mathrm{task}}$ \\
\midrule
\texttt{suicide\_risk}           & 40 & 0.250 & 0.350 & 0.100 & 110.3 \\
\texttt{what\_is\_the\_tao}      & 36 & 0.500 & 0.639 & 0.139 & 70.8 \\
\texttt{logical\_sequence}       & 39 & 0.252 & 0.385 & 0.133 & 50.1 \\
\texttt{mathematical\_induction} & 70 & 0.500 & 0.609 & 0.109 & 38.6 \\
\texttt{empirical\_judgments}    & 99 & 0.333 & 0.455 & 0.121 & 8.6 \\
\texttt{sentence\_ambiguity}     & 60 & 0.500 & 0.667 & 0.167 & 7.2 \\
\texttt{common\_morpheme}        & 50 & 0.250 & 0.420 & 0.170 & 6.7 \\
\bottomrule
\end{tabular}
\end{table}

\section*{Declaration of competing interest}

The author declares no competing interests.

\section*{CRediT author statement}

Zeyu Fu: Conceptualization, Methodology, Software, Investigation, Formal
analysis, Writing -- original draft, Writing -- review \& editing.

\section*{Funding}
This research received no external funding.

\section*{Data and code availability}
All experiments are synthetic. The complete training and analysis code, the per-cell
logs for the three negative grids, and the calibration/figure-generation scripts required
to reproduce every reported result and figure are archived in a public repository under
open licenses (code: MIT; data and figures: CC~BY~4.0). Repository:
\href{https://github.com/PeterPonyu/calibration-traps}{\nolinkurl{github.com/PeterPonyu/calibration-traps}};
archived snapshot DOI:
\href{https://doi.org/10.5281/zenodo.21020386}{\texttt{10.5281/zenodo.21020386}}.

\bibliographystyle{plainnat}
\section*{Declaration of generative AI and AI-assisted technologies}

During the preparation of this work the author used Claude (Anthropic) to assist
with drafting and editing the manuscript text, developing and reviewing the
analysis code, and preparing the figures. After using this tool, the author
reviewed and edited the content as needed and takes full responsibility for the
content of the published article.

\begin{thebibliography}{21}
\providecommand{\natexlab}[1]{#1}
\providecommand{\url}[1]{\texttt{#1}}
\providecommand{\doi}[1]{\href{https://doi.org/#1}{\path{doi:#1}}}

\bibitem[Agarwal et~al.(2021)Agarwal, Schwarzer, Castro, Courville, and
  Bellemare]{agarwal2021precipice}
Rishabh Agarwal, Max Schwarzer, Pablo~Samuel Castro, Aaron Courville, and
  Marc~G. Bellemare.
\newblock Deep reinforcement learning at the edge of the statistical
  precipice.
\newblock In \emph{Advances in Neural Information Processing Systems
  (NeurIPS)}, 2021. arXiv:2108.13264.
\newblock \doi{10.48550/arXiv.2108.13264}
\bibitem[Aky{\"u}rek et~al.(2024)Aky{\"u}rek, Wang, Kim, and
  Andreas]{akyurek2024iclalgorithms}
Ekin Aky{\"u}rek, Bailin Wang, Yoon Kim, and Jacob Andreas.
\newblock In-context language learning: Architectures and algorithms.
\newblock In \emph{International Conference on Machine Learning (ICML)}, 2024.
  arXiv:2401.12973.
\newblock \doi{10.48550/arXiv.2401.12973}
\bibitem[Barak et~al.(2022)Barak, Edelman, Goel, Kakade, Malach, and
  Zhang]{barak2022hidden}
Boaz Barak, Benjamin~L. Edelman, Surbhi Goel, Sham Kakade, Eran Malach, and
  Cyril Zhang.
\newblock Hidden progress in deep learning: {SGD} learns parities near the
  computational limit.
\newblock In \emph{Advances in Neural Information Processing Systems
  (NeurIPS)}, 2022.
\newblock arXiv:2207.08799.
\newblock \doi{10.48550/arXiv.2207.08799}
\bibitem[Bengio et~al.(2009)Bengio, Louradour, Collobert, and
  Weston]{bengio2009curriculum}
Yoshua Bengio, J{\'e}r{\^o}me Louradour, Ronan Collobert, and Jason Weston.
\newblock Curriculum learning.
\newblock In \emph{Proceedings of the 26th International Conference on Machine
  Learning (ICML)}, pages 41--48, 2009.

\bibitem[Biderman et~al.(2023)Biderman, Schoelkopf, Anthony, Bradley, O'Brien,
  Hallahan, Khan, Purohit, Prashanth, Raff, Skowron, Sutawika, and van~der
  Wal]{biderman2023pythia}
Stella Biderman, Hailey Schoelkopf, Quentin Anthony, Herbie Bradley, Kyle
  O'Brien, Eric Hallahan, Mohammad~Aflah Khan, Shivanshu Purohit, USVSN~Sai
  Prashanth, Edward Raff, Aviya Skowron, Lintang Sutawika, and Oskar van~der
  Wal.
\newblock Pythia: A suite for analyzing large language models across training
  and scaling.
\newblock In \emph{International Conference on Machine Learning (ICML)}, 2023.
  arXiv:2304.01373.
\newblock \doi{10.48550/arXiv.2304.01373}
\bibitem[Cohen et~al.(2021)Cohen, Kaur, Li, Kolter, and
  Talwalkar]{cohen2021eos}
Jeremy~M. Cohen, Simran Kaur, Yuanzhi Li, J.~Zico Kolter, and Ameet Talwalkar.
\newblock Gradient descent on neural networks typically occurs at the edge of
  stability.
\newblock In \emph{International Conference on Learning Representations
  (ICLR)}, 2021.
\newblock arXiv:2103.00065.
\newblock \doi{10.48550/arXiv.2103.00065}
\bibitem[Du et~al.(2024)Du, Zeng, Dong, and Tang]{du2024loss}
Zhengxiao Du, Aohan Zeng, Yuxiao Dong, and Jie Tang.
\newblock Understanding emergent abilities of language models from the loss
  perspective.
\newblock In \emph{Advances in Neural Information Processing Systems
  (NeurIPS)}, 2024. arXiv:2403.15796.
\newblock \doi{10.48550/arXiv.2403.15796}
\bibitem[Edelman et~al.(2024)Edelman, Edelman, Goel, Malach, and
  Tsilivis]{edelman2024statinduction}
Benjamin~L. Edelman, Ezra Edelman, Surbhi Goel, Eran Malach, and Nikolaos
  Tsilivis.
\newblock The evolution of statistical induction heads: In-context learning
  {Markov} chains.
\newblock In \emph{Advances in Neural Information Processing Systems
  (NeurIPS)}, 2024. arXiv:2402.11004.
\newblock \doi{10.48550/arXiv.2402.11004}
\bibitem[Feller(1968)]{feller1968}
William Feller.
\newblock \emph{An Introduction to Probability Theory and Its Applications},
  volume~1.
\newblock Wiley, New York, 3rd edition, 1968.

\bibitem[Fu and Lou(2003)]{fulou2003runs}
James~C. Fu and W.~Y. Wendy Lou.
\newblock \emph{Distribution Theory of Runs and Patterns and Its Applications: A
  Finite Markov Chain Imbedding Approach}.
\newblock World Scientific, Singapore, 2003.

\bibitem[Garg et~al.(2022)Garg, Tsipras, Liang, and Valiant]{garg2022icl}
Shivam Garg, Dimitris Tsipras, Percy Liang, and Gregory Valiant.
\newblock What can transformers learn in-context? {A} case study of simple
  function classes.
\newblock In \emph{Advances in Neural Information Processing Systems
  (NeurIPS)}, 2022. arXiv:2208.01066.
\newblock \doi{10.48550/arXiv.2208.01066}
\bibitem[Glaz et~al.(2001)Glaz, Naus, and Wallenstein]{glaz2001scan}
Joseph Glaz, Joseph Naus, and Sylvan Wallenstein.
\newblock \emph{Scan Statistics}.
\newblock Springer, New York, 2001.

\bibitem[Kumar et~al.(2024)Kumar, Bordelon, Gershman, and
  Pehlevan]{kumar2024lazytorich}
Tanishq Kumar, Blake Bordelon, Samuel~J. Gershman, and Cengiz Pehlevan.
\newblock Grokking as the transition from lazy to rich training dynamics.
\newblock In \emph{International Conference on Learning Representations
  (ICLR)}, 2024. arXiv:2310.06110.
\newblock \doi{10.48550/arXiv.2310.06110}
\bibitem[Laskin et~al.(2023)Laskin, Wang, Oh, Parisotto, Spencer, Steigerwald,
  Strouse, Hansen, Filos, Brooks, Gazeau, Sahni, Singh, and Mnih]{laskin2022ad}
Michael Laskin, Luyu Wang, Junhyuk Oh, Emilio Parisotto, Stephen Spencer, Richie
  Steigerwald, DJ~Strouse, Steven Hansen, Angelos Filos, Ethan Brooks, Maxime
  Gazeau, Himanshu Sahni, Satinder Singh, and Volodymyr Mnih.
\newblock In-context reinforcement learning with algorithm distillation.
\newblock In \emph{International Conference on Learning Representations (ICLR)}, 2023.
  arXiv:2210.14215.
\newblock \doi{10.48550/arXiv.2210.14215}
\bibitem[Lee et~al.(2023)Lee, Xie, Pacchiano, Chandak, Finn, Nachum, and
  Brunskill]{lee2023icrl}
Jonathan~N. Lee, Annie Xie, Aldo Pacchiano, Yash Chandak, Chelsea Finn, Ofir
  Nachum, and Emma Brunskill.
\newblock Supervised pretraining can learn in-context reinforcement learning.
\newblock In \emph{Advances in Neural Information Processing Systems
  (NeurIPS)}, 2023.
\newblock arXiv:2306.14892.
\newblock \doi{10.48550/arXiv.2306.14892}
\bibitem[Liu et~al.(2023)Liu, Michaud, and Tegmark]{liu2023omnigrok}
Ziming Liu, Eric~J. Michaud, and Max Tegmark.
\newblock Omnigrok: Grokking beyond algorithmic data.
\newblock In \emph{International Conference on Learning Representations
  (ICLR)}, 2023.
\newblock arXiv:2210.01117.
\newblock \doi{10.48550/arXiv.2210.01117}
\bibitem[McKenzie et~al.(2023)McKenzie, Lyzhov, Pieler, Parrish, Mueller,
  Prabhu, McLean, Kirtland, Ross, Liu, Gritsevskiy, Wurgaft, Kauffman, Recchia,
  Liu, Cavanagh, Weiss, Huang, Tseng, Korbak, Shen, Zhang, Zhou, Kim, Bowman,
  and Perez]{mckenzie2023inversescaling}
Ian~R. McKenzie, Alexander Lyzhov, Michael Pieler, Alicia Parrish, Aaron
  Mueller, Ameya Prabhu, Euan McLean, Aaron Kirtland, Alexis Ross, Alisa Liu,
  Andrew Gritsevskiy, Daniel Wurgaft, Derik Kauffman, Gabriel Recchia, Jiacheng
  Liu, Joe Cavanagh, Max Weiss, Sicong Huang, {The Floating Droid}, Tom Tseng, Tomasz Korbak, Xudong
  Shen, Yuhui Zhang, Zhengping Zhou, Najoung Kim, Samuel~R. Bowman, and Ethan
  Perez.
\newblock Inverse scaling: When bigger isn't better.
\newblock \emph{Transactions on Machine Learning Research (TMLR)}, 2023.
\newblock arXiv:2306.09479.
\newblock \doi{10.48550/arXiv.2306.09479}
\bibitem[Nanda et~al.(2023)Nanda, Chan, Lieberum, Smith, and
  Steinhardt]{nanda2023progress}
Neel Nanda, Lawrence Chan, Tom Lieberum, Jess Smith, and Jacob Steinhardt.
\newblock Progress measures for grokking via mechanistic interpretability.
\newblock In \emph{International Conference on Learning Representations
  (ICLR)}, 2023.
\newblock arXiv:2301.05217.
\newblock \doi{10.48550/arXiv.2301.05217}
\bibitem[Olsson et~al.(2022)Olsson, Elhage, Nanda, Joseph, DasSarma, Henighan,
  Mann, Askell, Bai, Chen, Conerly, Drain, Ganguli, {Hatfield-Dodds},
  Hernandez, Johnston, Jones, Kernion, Lovitt, Ndousse, Amodei, Brown, Clark,
  Kaplan, McCandlish, and Olah]{olsson2022induction}
Catherine Olsson, Nelson Elhage, Neel Nanda, Nicholas Joseph, Nova DasSarma,
  Tom Henighan, Ben Mann, Amanda Askell, Yuntao Bai, Anna Chen, Tom Conerly,
  Dawn Drain, Deep Ganguli, Zac {Hatfield-Dodds}, Danny Hernandez, Scott
  Johnston, Andy Jones, Jackson Kernion, Liane Lovitt, Kamal Ndousse, Dario
  Amodei, Tom Brown, Jack Clark, Jared Kaplan, Sam McCandlish, and Chris Olah.
\newblock In-context learning and induction heads.
\newblock \emph{arXiv preprint} arXiv:2209.11895, 2022.
\newblock \doi{10.48550/arXiv.2209.11895}
\bibitem[Oneto et~al.(2026)Oneto, Ridella, Minisi, Coraddu, and Anguita]{oneto2026}
Luca Oneto, Sandro Ridella, Simone Minisi, Andrea Coraddu, and Davide Anguita.
\newblock Reconciling grokking with statistical learning theory through the lens
  of norm- and stability-based generalization bounds.
\newblock \emph{Neurocomputing}, 674:\penalty0 132826, 2026.
\newblock \doi{10.1016/j.neucom.2026.132826}
\bibitem[Paperno et~al.(2016)Paperno, Kruszewski, Lazaridou, Pham, Bernardi,
  Pezzelle, Baroni, Boleda, and Fern{\'a}ndez]{paperno2016lambada}
Denis Paperno, Germ{\'a}n Kruszewski, Angeliki Lazaridou, Quan~Ngoc Pham,
  Raffaella Bernardi, Sandro Pezzelle, Marco Baroni, Gemma Boleda, and Raquel
  Fern{\'a}ndez.
\newblock The {LAMBADA} dataset: Word prediction requiring a broad discourse
  context.
\newblock In \emph{Proceedings of the 54th Annual Meeting of the Association
  for Computational Linguistics (ACL)}, pages 1525--1534, 2016.
  arXiv:1606.06031.
\newblock \doi{10.18653/v1/P16-1144}
\bibitem[Pineau et~al.(2021)Pineau, Vincent-Lamarre, Sinha, Larivi{\`e}re,
  Beygelzimer, {d'Alch{\'e}-Buc}, Fox, and
  Larochelle]{pineau2021reproducibility}
Joelle Pineau, Philippe Vincent-Lamarre, Koustuv Sinha, Vincent Larivi{\`e}re,
  Alina Beygelzimer, Florence {d'Alch{\'e}-Buc}, Emily Fox, and Hugo
  Larochelle.
\newblock Improving reproducibility in machine learning research ({A} report
  from the {NeurIPS} 2019 reproducibility program).
\newblock \emph{Journal of Machine Learning Research}, 22\penalty0
  (164):\penalty0 1--20, 2021.
\newblock arXiv:2003.12206.
\newblock \doi{10.48550/arXiv.2003.12206}
\bibitem[Power et~al.(2022)Power, Burda, Edwards, Babuschkin, and
  Misra]{power2022grokking}
Alethea Power, Yuri Burda, Harri Edwards, Igor Babuschkin, and Vedant Misra.
\newblock Grokking: Generalization beyond overfitting on small algorithmic
  datasets.
\newblock \emph{arXiv preprint} arXiv:2201.02177, 2022.
\newblock \doi{10.48550/arXiv.2201.02177}
\bibitem[Prieto et~al.(2025)Prieto, Barsbey, Mediano, and
  Birdal]{prieto2025stablemax}
Lucas Prieto, Melih Barsbey, Pedro A.~M. Mediano, and Tolga Birdal.
\newblock Grokking at the edge of numerical stability.
\newblock In \emph{International Conference on Learning Representations
  (ICLR)}, 2025.
\newblock arXiv:2501.04697.
\newblock \doi{10.48550/arXiv.2501.04697}
\bibitem[Sabry and Belz(2025)]{sabry2025induction}
Mohammed Sabry and Anya Belz.
\newblock Induction signatures are not enough: A matched-compute study of
  load-bearing structure in in-context learning.
\newblock arXiv:2509.22947, 2025.
\newblock \doi{10.48550/arXiv.2509.22947}
\bibitem[Schaeffer et~al.(2023)Schaeffer, Miranda, and
  Koyejo]{schaeffer2023mirage}
Rylan Schaeffer, Brando Miranda, and Sanmi Koyejo.
\newblock Are emergent abilities of large language models a mirage?
\newblock In \emph{Advances in Neural Information Processing Systems
  (NeurIPS)}, 2023.
\newblock arXiv:2304.15004.
\newblock \doi{10.48550/arXiv.2304.15004}
\bibitem[Snell et~al.(2024)Snell, Wallace, Klein, and Levine]{snell2024predicting}
Charlie Snell, Eric Wallace, Dan Klein, and Sergey Levine.
\newblock Predicting emergent capabilities by finetuning.
\newblock arXiv:2411.16035, 2024.
\newblock \doi{10.48550/arXiv.2411.16035}
\bibitem[Soviany et~al.(2022)Soviany, Ionescu, Rota, and
  Sebe]{soviany2022curriculum}
Petru Soviany, Radu~Tudor Ionescu, Paolo Rota, and Nicu Sebe.
\newblock Curriculum learning: A survey.
\newblock \emph{International Journal of Computer Vision}, 130:\penalty0
  1526--1565, 2022.
\newblock \doi{10.1007/s11263-022-01611-x}

\bibitem[Srivastava et~al.(2022)]{srivastava2022bigbench}
Aarohi Srivastava et~al.
\newblock Beyond the imitation game: Quantifying and extrapolating the
  capabilities of language models.
\newblock \emph{Transactions on Machine Learning Research}, 2022.
  arXiv:2206.04615.
\newblock \doi{10.48550/arXiv.2206.04615}

\bibitem[Tikeng Notsawo et~al.(2023)Tikeng Notsawo, Zhou, Pezeshki, Rish, and
  Dumas]{notsawo2023predicting}
Pascal~Jr. {Tikeng Notsawo}, Hattie Zhou, Mohammad Pezeshki, Irina Rish, and Guillaume
  Dumas.
\newblock Predicting grokking long before it happens: A look into the loss
  landscape of models which grok.
\newblock \emph{arXiv preprint} arXiv:2306.13253, 2023.
\newblock \doi{10.48550/arXiv.2306.13253}
\bibitem[von Oswald et~al.(2023)von Oswald, Niklasson, Randazzo, Sacramento,
  Mordvintsev, Zhmoginov, and Vladymyrov]{vonoswald2023gd}
Johannes von Oswald, Eyvind Niklasson, Ettore Randazzo, Jo{\~a}o Sacramento,
  Alexander Mordvintsev, Andrey Zhmoginov, and Max Vladymyrov.
\newblock Transformers learn in-context by gradient descent.
\newblock In \emph{International Conference on Machine Learning (ICML)},
  2023. arXiv:2212.07677.
\newblock \doi{10.48550/arXiv.2212.07677}
\bibitem[Wang et~al.(2024)Wang, Blaser, Daneshmand, and
  Zhang]{wang2024tdtransformers}
Jiuqi Wang, Ethan Blaser, Hadi Daneshmand, and Shangtong Zhang.
\newblock Transformers can learn temporal difference methods for in-context
  reinforcement learning.
\newblock arXiv:2405.13861, 2024.
\newblock \doi{10.48550/arXiv.2405.13861}
\bibitem[Wei et~al.(2022)Wei, Tay, Bommasani, Raffel, Zoph, Borgeaud, Yogatama,
  Bosma, Zhou, Metzler, Chi, Hashimoto, Vinyals, Liang, Dean, and
  Fedus]{wei2022emergent}
Jason Wei, Yi~Tay, Rishi Bommasani, Colin Raffel, Barret Zoph, Sebastian
  Borgeaud, Dani Yogatama, Maarten Bosma, Denny Zhou, Donald Metzler, Ed~H.
  Chi, Tatsunori Hashimoto, Oriol Vinyals, Percy Liang, Jeff Dean, and William
  Fedus.
\newblock Emergent abilities of large language models.
\newblock \emph{Transactions on Machine Learning Research (TMLR)}, 2022.
\newblock arXiv:2206.07682.
\newblock \doi{10.48550/arXiv.2206.07682}
\bibitem[Wu et~al.(2021)Wu, Dyer, and Neyshabur]{wu2021curricula}
Xiaoxia Wu, Ethan Dyer, and Behnam Neyshabur.
\newblock When do curricula work?
\newblock In \emph{International Conference on Learning Representations
  (ICLR)}, 2021. arXiv:2012.03107.
\newblock \doi{10.48550/arXiv.2012.03107}
\bibitem[Xu(2026)]{xu2026earlywarning}
Yongzhong Xu.
\newblock Early-warning signals of grokking via loss-landscape geometry.
\newblock \emph{arXiv preprint} arXiv:2602.16967, 2026.
\newblock \doi{10.48550/arXiv.2602.16967}
\bibitem[Zisman et~al.(2024)Zisman, Kurenkov, Nikulin, Sinii, and
  Kolesnikov]{zisman2024noise}
Ilya Zisman, Vladislav Kurenkov, Alexander Nikulin, Viacheslav Sinii, and Sergey
  Kolesnikov.
\newblock Emergence of in-context reinforcement learning from noise
  distillation.
\newblock In \emph{International Conference on Machine Learning (ICML)}, 2024.
\newblock arXiv:2312.12275.
\newblock \doi{10.48550/arXiv.2312.12275}
\end{thebibliography}

\end{document}
